% Figure: staggered control volumes and primal--dual fluxes (Chapter 6).
% Cells, solid, staggered control volumes, primal faces, dual faces and the
% interpolation weights read from the code's schedules:
% generated/momentum_dual.tex (momentum.jl: uniform, composite and disc
% fixtures). Labels hand-written.
\begin{figure}
  \centering
  \def\fxmnone{%
    \def\fxmcell##1##2##3##4##5##6##7{}\def\fxmsolid##1##2{}\def\fxmwall##1##2##3##4##5{}%
    \def\fxmdual##1##2##3##4##5##6{}\def\fxmopen##1##2{}\def\fxmface##1##2##3##4##5##6##7{}%
    \def\fxmlink##1##2##3##4##5##6##7##8##9{}\def\fxmdualface##1##2##3##4##5##6##7{}%
    \def\fxmhanging##1##2##3{}}%
  % Passes of panel #1: 1 level shading, open strips and highlighted control
  % volume, 2 solid, 3 cells and control volumes, 4 primal faces.
  \def\fxmpasses#1{%
    \fxmnone
    \def\fxmcell##1##2##3##4##5##6##7{\ifnum##1=#1 \fill[tr level ##6] (##2,##3) rectangle (##4,##5);\fi}%
    \def\fxmopen##1##2{\ifnum##1=#1 \fill[tr dual fill] ##2;\fi}%
    \def\fxmdual##1##2##3##4##5##6{\ifnum##1=#1 \ifnum##6=1 \ifnum#1<3
      \fill[tr dual fill] (##2,##3) rectangle (##4,##5);\fi\fi\fi}%
    % Generated by docs/techref/figures/generate.jl; do not edit.
\fxmcell{1}{4}{4}{5}{5}{0}{0}
\fxmcell{1}{5}{4}{6}{5}{0}{0}
\fxmcell{1}{6}{4}{7}{5}{0}{0}
\fxmcell{1}{7}{4}{8}{5}{0}{0}
\fxmcell{1}{4}{5}{5}{6}{0}{0}
\fxmcell{1}{5}{5}{6}{6}{0}{0}
\fxmcell{1}{6}{5}{7}{6}{0}{0}
\fxmcell{1}{7}{5}{8}{6}{0}{0}
\fxmcell{1}{4}{6}{5}{7}{0}{0}
\fxmcell{1}{5}{6}{6}{7}{0}{0}
\fxmcell{1}{6}{6}{7}{7}{0}{0}
\fxmcell{1}{7}{6}{8}{7}{0}{0}
\fxmdual{1}{4.5}{4}{5.5}{5}{0}
\fxmdual{1}{4.5}{5}{5.5}{6}{0}
\fxmdual{1}{4.5}{6}{5.5}{7}{0}
\fxmdual{1}{5.5}{4}{6.5}{5}{0}
\fxmdual{1}{5.5}{5}{6.5}{6}{1}
\fxmdual{1}{5.5}{6}{6.5}{7}{0}
\fxmdual{1}{6.5}{4}{7.5}{5}{0}
\fxmdual{1}{6.5}{5}{7.5}{6}{0}
\fxmdual{1}{6.5}{6}{7.5}{7}{0}
\fxmface{1}{1}{6}{5.5}{1}{f}{1.00}
\fxmface{1}{2}{6.5}{6}{1}{k'}{1.00}
\fxmface{1}{2}{5.5}{5}{1}{l}{1.00}
\fxmface{1}{2}{6.5}{5}{1}{l'}{1.00}
\fxmface{1}{1}{7}{5.5}{1}{g'}{1.00}
\fxmface{1}{1}{5}{5.5}{1}{g}{1.00}
\fxmface{1}{2}{5.5}{6}{1}{k}{1.00}
\fxmdualface{1}{5.5}{5}{5.5}{6}{\tfrac{1}{2}\,f+\tfrac{1}{2}\,g}{l}
\fxmdualface{1}{6.5}{5}{6.5}{6}{\tfrac{1}{2}\,f+\tfrac{1}{2}\,g'}{r}
\fxmdualface{1}{5.5}{5}{6.5}{5}{\tfrac{1}{2}\,l+\tfrac{1}{2}\,l'}{b}
\fxmdualface{1}{5.5}{6}{6.5}{6}{\tfrac{1}{2}\,k+\tfrac{1}{2}\,k'}{t}
\fxmcell{2}{2}{3}{3}{4}{0}{0}
\fxmcell{2}{3}{3}{4}{4}{0}{0}
\fxmcell{2}{4}{3}{5}{4}{0}{0}
\fxmcell{2}{5}{3}{6}{4}{0}{0}
\fxmcell{2}{2}{4}{3}{5}{0}{0}
\fxmcell{2}{3}{4}{4}{5}{0}{0}
\fxmcell{2}{2}{5}{3}{6}{0}{0}
\fxmcell{2}{3}{5}{4}{6}{0}{0}
\fxmcell{2}{2}{6}{3}{7}{0}{0}
\fxmcell{2}{3}{6}{4}{7}{0}{0}
\fxmcell{2}{2}{7}{3}{8}{0}{0}
\fxmcell{2}{3}{7}{4}{8}{0}{0}
\fxmcell{2}{4}{4}{4.5}{4.5}{1}{0}
\fxmcell{2}{4.5}{4}{5}{4.5}{1}{0}
\fxmcell{2}{5}{4}{5.5}{4.5}{1}{0}
\fxmcell{2}{4}{4.5}{4.5}{5}{1}{0}
\fxmcell{2}{4.5}{4.5}{5}{5}{1}{0}
\fxmcell{2}{5}{4.5}{5.5}{5}{1}{0}
\fxmcell{2}{4}{5}{4.5}{5.5}{1}{0}
\fxmcell{2}{4.5}{5}{5}{5.5}{1}{0}
\fxmcell{2}{5}{5}{5.5}{5.5}{1}{0}
\fxmcell{2}{4}{5.5}{4.5}{6}{1}{0}
\fxmcell{2}{4.5}{5.5}{5}{6}{1}{0}
\fxmcell{2}{5}{5.5}{5.5}{6}{1}{0}
\fxmcell{2}{4}{6}{4.5}{6.5}{1}{0}
\fxmcell{2}{4.5}{6}{5}{6.5}{1}{0}
\fxmcell{2}{5}{6}{5.5}{6.5}{1}{0}
\fxmcell{2}{4}{6.5}{4.5}{7}{1}{0}
\fxmcell{2}{4.5}{6.5}{5}{7}{1}{0}
\fxmcell{2}{5}{6.5}{5.5}{7}{1}{0}
\fxmcell{2}{4}{7}{4.5}{7.5}{1}{0}
\fxmcell{2}{4.5}{7}{5}{7.5}{1}{0}
\fxmcell{2}{5}{7}{5.5}{7.5}{1}{0}
\fxmdual{2}{1.5}{3}{2.5}{4}{0}
\fxmdual{2}{1.5}{4}{2.5}{5}{0}
\fxmdual{2}{1.5}{5}{2.5}{6}{0}
\fxmdual{2}{1.5}{6}{2.5}{7}{0}
\fxmdual{2}{1.5}{7}{2.5}{8}{0}
\fxmdual{2}{2.5}{3}{3.5}{4}{0}
\fxmdual{2}{2.5}{4}{3.5}{5}{0}
\fxmdual{2}{2.5}{5}{3.5}{6}{0}
\fxmdual{2}{2.5}{6}{3.5}{7}{0}
\fxmdual{2}{2.5}{7}{3.5}{8}{0}
\fxmdual{2}{3.5}{3}{4.5}{4}{0}
\fxmdual{2}{3.5}{4}{4.25}{4.5}{0}
\fxmdual{2}{3.5}{4.5}{4.25}{5}{0}
\fxmdual{2}{3.5}{5}{4.25}{5.5}{1}
\fxmdual{2}{3.5}{5.5}{4.25}{6}{0}
\fxmdual{2}{3.5}{6}{4.25}{6.5}{0}
\fxmdual{2}{3.5}{6.5}{4.25}{7}{0}
\fxmdual{2}{3.5}{7}{4.25}{7.5}{0}
\fxmdual{2}{4.25}{4}{4.75}{4.5}{0}
\fxmdual{2}{4.25}{4.5}{4.75}{5}{0}
\fxmdual{2}{4.25}{5}{4.75}{5.5}{0}
\fxmdual{2}{4.25}{5.5}{4.75}{6}{0}
\fxmdual{2}{4.25}{6}{4.75}{6.5}{0}
\fxmdual{2}{4.25}{6.5}{4.75}{7}{0}
\fxmdual{2}{4.25}{7}{4.75}{7.5}{0}
\fxmdual{2}{4.5}{3}{5.5}{4}{0}
\fxmdual{2}{4.75}{4}{5.25}{4.5}{0}
\fxmdual{2}{4.75}{4.5}{5.25}{5}{0}
\fxmdual{2}{4.75}{5}{5.25}{5.5}{0}
\fxmdual{2}{4.75}{5.5}{5.25}{6}{0}
\fxmdual{2}{4.75}{6}{5.25}{6.5}{0}
\fxmdual{2}{4.75}{6.5}{5.25}{7}{0}
\fxmdual{2}{4.75}{7}{5.25}{7.5}{0}
\fxmface{2}{1}{4}{5.75}{0.5}{f_2}{1.00}
\fxmface{2}{1}{4}{5.25}{0.5}{f_1}{1.00}
\fxmface{2}{1}{4.5}{5.25}{0.5}{h_1}{1.00}
\fxmface{2}{2}{3.5}{6}{1}{\bar{k}^{+}}{1.00}
\fxmface{2}{2}{4.25}{5.5}{0.5}{k}{1.00}
\fxmface{2}{1}{3}{5.5}{1}{g}{1.00}
\fxmface{2}{2}{3.5}{5}{1}{\bar{k}^{-}}{1.00}
\fxmface{2}{2}{4.25}{5}{0.5}{k_1}{1.00}
\fxmlink{2}{4}{5.25}{3}{5.5}{1}{1}{\tfrac{1}{2}\,f_1+\tfrac{1}{4}\,g}{0}
\fxmlink{2}{4}{5.25}{4}{4.25}{2}{1}{\tfrac{1}{8}\,\bar{k}^{-}}{-1}
\fxmlink{2}{4}{5.25}{4}{4.75}{2}{1}{\tfrac{1}{8}\,\bar{k}^{-}+\tfrac{1}{2}\,k_1}{-1}
\fxmlink{2}{4}{5.25}{4}{5.75}{1}{1}{-\tfrac{1}{4}\,f_1+\tfrac{1}{4}\,f_2}{1}
\fxmlink{2}{4}{5.25}{4}{5.75}{2}{1}{\tfrac{1}{2}\,k}{-1}
\fxmlink{2}{4}{5.25}{4}{6.25}{2}{1}{\tfrac{1}{8}\,\bar{k}^{+}}{-1}
\fxmlink{2}{4}{5.25}{4}{6.75}{2}{1}{\tfrac{1}{8}\,\bar{k}^{+}}{-1}
\fxmlink{2}{4}{5.25}{4.5}{5.25}{1}{1}{\tfrac{1}{2}\,f_1+\tfrac{1}{2}\,h_1}{0}
\fxmhanging{2}{4}{4.5}
\fxmhanging{2}{4}{4.5}
\fxmhanging{2}{4}{5.5}
\fxmhanging{2}{4}{5.5}
\fxmhanging{2}{4}{6.5}
\fxmhanging{2}{4}{6.5}
\fxmcell{3}{6}{4.5}{6.5}{5}{1}{0}
\fxmcell{3}{6.5}{4.5}{7}{5}{1}{0}
\fxmcell{3}{7}{4.5}{7.5}{5}{1}{1}
\fxmsolid{3}{(7.5,4.894926) -- (7.5,5) -- (7.244029,5) -- cycle}
\fxmwall{3}{7.5}{4.894926}{7.244029}{5}
\fxmcell{3}{7.5}{4.5}{8}{5}{1}{1}
\fxmsolid{3}{(8,4.783878) -- (8,5) -- (7.5,5) -- (7.5,4.857327) -- cycle}
\fxmwall{3}{8}{4.783878}{7.5}{4.857327}
\fxmcell{3}{6}{5}{6.5}{5.5}{1}{0}
\fxmcell{3}{6.5}{5}{7}{5.5}{1}{1}
\fxmsolid{3}{(7,5.165022) -- (7,5.5) -- (6.774767,5.5) -- cycle}
\fxmwall{3}{7}{5.165022}{6.774767}{5.5}
\fxmcell{3}{7}{5}{7.5}{5.5}{1}{1}
\fxmsolid{3}{(7.22103,5) -- (7.5,5) -- (7.5,5.5) -- (7,5.5) -- (7,5.238627) -- cycle}
\fxmwall{3}{7.22103}{5}{7}{5.238627}
\fxmcell{3}{7.5}{5}{8}{5.5}{1}{2}
\fxmsolid{3}{(7.5,5) -- (8,5) -- (8,5.5) -- (7.5,5.5) -- cycle}
\fxmcell{3}{6}{5.5}{6.5}{6}{1}{0}
\fxmcell{3}{6.5}{5.5}{7}{6}{1}{1}
\fxmsolid{3}{(6.753218,5.5) -- (7,5.5) -- (7,6) -- (6.528991,6) -- cycle}
\fxmwall{3}{6.753218}{5.5}{6.528991}{6}
\fxmcell{3}{7}{5.5}{7.5}{6}{1}{2}
\fxmsolid{3}{(7,5.5) -- (7.5,5.5) -- (7.5,6) -- (7,6) -- cycle}
\fxmcell{3}{7.5}{5.5}{8}{6}{1}{2}
\fxmsolid{3}{(7.5,5.5) -- (8,5.5) -- (8,6) -- (7.5,6) -- cycle}
\fxmopen{3}{(6.5,5) -- (7,5) -- (7,5.165022) -- (6.774767,5.5) -- (6.5,5.5) -- cycle}
\fxmopen{3}{(7,5) -- (7.22103,5) -- (7,5.238627) -- cycle}
\fxmdual{3}{6.75}{5}{7.25}{5.5}{1}
\fxmface{3}{2}{7.25}{5}{0.5}{l'}{0.47}
\fxmface{3}{2}{6.75}{5.5}{0.5}{k}{0.53}
\fxmface{3}{2}{7.25}{5.5}{0.5}{k'}{0.00}
\fxmface{3}{1}{7.5}{5.25}{0.5}{g'}{0.00}
\fxmface{3}{1}{6.5}{5.25}{0.5}{g}{1.00}
\fxmface{3}{1}{7}{5.25}{0.5}{f}{0.40}
\fxmface{3}{2}{6.75}{5}{0.5}{l}{1.00}
\fxmdualface{3}{6.75}{5}{6.75}{5.5}{\tfrac12(\aperture{f}\,f+\aperture{g}\,g)}{l}
\fxmdualface{3}{6.75}{5}{7.25}{5}{\tfrac12(\aperture{l}\,l+\aperture{l'}\,l')}{b}
\fxmdualface{3}{6.75}{5.5}{7.25}{5.5}{\tfrac12(\aperture{k}\,k+\aperture{k'}\,k')}{t}
\fxmdualface{3}{7.25}{5}{7.25}{5.5}{\tfrac12(\aperture{f}\,f+\aperture{g'}\,g')}{r}
%
    \fxmnone
    \def\fxmsolid##1##2{\ifnum##1=#1 \fill[white] ##2; \fill[tr solid] ##2;\fi}%
    % Generated by docs/techref/figures/generate.jl; do not edit.
\fxmcell{1}{4}{4}{5}{5}{0}{0}
\fxmcell{1}{5}{4}{6}{5}{0}{0}
\fxmcell{1}{6}{4}{7}{5}{0}{0}
\fxmcell{1}{7}{4}{8}{5}{0}{0}
\fxmcell{1}{4}{5}{5}{6}{0}{0}
\fxmcell{1}{5}{5}{6}{6}{0}{0}
\fxmcell{1}{6}{5}{7}{6}{0}{0}
\fxmcell{1}{7}{5}{8}{6}{0}{0}
\fxmcell{1}{4}{6}{5}{7}{0}{0}
\fxmcell{1}{5}{6}{6}{7}{0}{0}
\fxmcell{1}{6}{6}{7}{7}{0}{0}
\fxmcell{1}{7}{6}{8}{7}{0}{0}
\fxmdual{1}{4.5}{4}{5.5}{5}{0}
\fxmdual{1}{4.5}{5}{5.5}{6}{0}
\fxmdual{1}{4.5}{6}{5.5}{7}{0}
\fxmdual{1}{5.5}{4}{6.5}{5}{0}
\fxmdual{1}{5.5}{5}{6.5}{6}{1}
\fxmdual{1}{5.5}{6}{6.5}{7}{0}
\fxmdual{1}{6.5}{4}{7.5}{5}{0}
\fxmdual{1}{6.5}{5}{7.5}{6}{0}
\fxmdual{1}{6.5}{6}{7.5}{7}{0}
\fxmface{1}{1}{6}{5.5}{1}{f}{1.00}
\fxmface{1}{2}{6.5}{6}{1}{k'}{1.00}
\fxmface{1}{2}{5.5}{5}{1}{l}{1.00}
\fxmface{1}{2}{6.5}{5}{1}{l'}{1.00}
\fxmface{1}{1}{7}{5.5}{1}{g'}{1.00}
\fxmface{1}{1}{5}{5.5}{1}{g}{1.00}
\fxmface{1}{2}{5.5}{6}{1}{k}{1.00}
\fxmdualface{1}{5.5}{5}{5.5}{6}{\tfrac{1}{2}\,f+\tfrac{1}{2}\,g}{l}
\fxmdualface{1}{6.5}{5}{6.5}{6}{\tfrac{1}{2}\,f+\tfrac{1}{2}\,g'}{r}
\fxmdualface{1}{5.5}{5}{6.5}{5}{\tfrac{1}{2}\,l+\tfrac{1}{2}\,l'}{b}
\fxmdualface{1}{5.5}{6}{6.5}{6}{\tfrac{1}{2}\,k+\tfrac{1}{2}\,k'}{t}
\fxmcell{2}{2}{3}{3}{4}{0}{0}
\fxmcell{2}{3}{3}{4}{4}{0}{0}
\fxmcell{2}{4}{3}{5}{4}{0}{0}
\fxmcell{2}{5}{3}{6}{4}{0}{0}
\fxmcell{2}{2}{4}{3}{5}{0}{0}
\fxmcell{2}{3}{4}{4}{5}{0}{0}
\fxmcell{2}{2}{5}{3}{6}{0}{0}
\fxmcell{2}{3}{5}{4}{6}{0}{0}
\fxmcell{2}{2}{6}{3}{7}{0}{0}
\fxmcell{2}{3}{6}{4}{7}{0}{0}
\fxmcell{2}{2}{7}{3}{8}{0}{0}
\fxmcell{2}{3}{7}{4}{8}{0}{0}
\fxmcell{2}{4}{4}{4.5}{4.5}{1}{0}
\fxmcell{2}{4.5}{4}{5}{4.5}{1}{0}
\fxmcell{2}{5}{4}{5.5}{4.5}{1}{0}
\fxmcell{2}{4}{4.5}{4.5}{5}{1}{0}
\fxmcell{2}{4.5}{4.5}{5}{5}{1}{0}
\fxmcell{2}{5}{4.5}{5.5}{5}{1}{0}
\fxmcell{2}{4}{5}{4.5}{5.5}{1}{0}
\fxmcell{2}{4.5}{5}{5}{5.5}{1}{0}
\fxmcell{2}{5}{5}{5.5}{5.5}{1}{0}
\fxmcell{2}{4}{5.5}{4.5}{6}{1}{0}
\fxmcell{2}{4.5}{5.5}{5}{6}{1}{0}
\fxmcell{2}{5}{5.5}{5.5}{6}{1}{0}
\fxmcell{2}{4}{6}{4.5}{6.5}{1}{0}
\fxmcell{2}{4.5}{6}{5}{6.5}{1}{0}
\fxmcell{2}{5}{6}{5.5}{6.5}{1}{0}
\fxmcell{2}{4}{6.5}{4.5}{7}{1}{0}
\fxmcell{2}{4.5}{6.5}{5}{7}{1}{0}
\fxmcell{2}{5}{6.5}{5.5}{7}{1}{0}
\fxmcell{2}{4}{7}{4.5}{7.5}{1}{0}
\fxmcell{2}{4.5}{7}{5}{7.5}{1}{0}
\fxmcell{2}{5}{7}{5.5}{7.5}{1}{0}
\fxmdual{2}{1.5}{3}{2.5}{4}{0}
\fxmdual{2}{1.5}{4}{2.5}{5}{0}
\fxmdual{2}{1.5}{5}{2.5}{6}{0}
\fxmdual{2}{1.5}{6}{2.5}{7}{0}
\fxmdual{2}{1.5}{7}{2.5}{8}{0}
\fxmdual{2}{2.5}{3}{3.5}{4}{0}
\fxmdual{2}{2.5}{4}{3.5}{5}{0}
\fxmdual{2}{2.5}{5}{3.5}{6}{0}
\fxmdual{2}{2.5}{6}{3.5}{7}{0}
\fxmdual{2}{2.5}{7}{3.5}{8}{0}
\fxmdual{2}{3.5}{3}{4.5}{4}{0}
\fxmdual{2}{3.5}{4}{4.25}{4.5}{0}
\fxmdual{2}{3.5}{4.5}{4.25}{5}{0}
\fxmdual{2}{3.5}{5}{4.25}{5.5}{1}
\fxmdual{2}{3.5}{5.5}{4.25}{6}{0}
\fxmdual{2}{3.5}{6}{4.25}{6.5}{0}
\fxmdual{2}{3.5}{6.5}{4.25}{7}{0}
\fxmdual{2}{3.5}{7}{4.25}{7.5}{0}
\fxmdual{2}{4.25}{4}{4.75}{4.5}{0}
\fxmdual{2}{4.25}{4.5}{4.75}{5}{0}
\fxmdual{2}{4.25}{5}{4.75}{5.5}{0}
\fxmdual{2}{4.25}{5.5}{4.75}{6}{0}
\fxmdual{2}{4.25}{6}{4.75}{6.5}{0}
\fxmdual{2}{4.25}{6.5}{4.75}{7}{0}
\fxmdual{2}{4.25}{7}{4.75}{7.5}{0}
\fxmdual{2}{4.5}{3}{5.5}{4}{0}
\fxmdual{2}{4.75}{4}{5.25}{4.5}{0}
\fxmdual{2}{4.75}{4.5}{5.25}{5}{0}
\fxmdual{2}{4.75}{5}{5.25}{5.5}{0}
\fxmdual{2}{4.75}{5.5}{5.25}{6}{0}
\fxmdual{2}{4.75}{6}{5.25}{6.5}{0}
\fxmdual{2}{4.75}{6.5}{5.25}{7}{0}
\fxmdual{2}{4.75}{7}{5.25}{7.5}{0}
\fxmface{2}{1}{4}{5.75}{0.5}{f_2}{1.00}
\fxmface{2}{1}{4}{5.25}{0.5}{f_1}{1.00}
\fxmface{2}{1}{4.5}{5.25}{0.5}{h_1}{1.00}
\fxmface{2}{2}{3.5}{6}{1}{\bar{k}^{+}}{1.00}
\fxmface{2}{2}{4.25}{5.5}{0.5}{k}{1.00}
\fxmface{2}{1}{3}{5.5}{1}{g}{1.00}
\fxmface{2}{2}{3.5}{5}{1}{\bar{k}^{-}}{1.00}
\fxmface{2}{2}{4.25}{5}{0.5}{k_1}{1.00}
\fxmlink{2}{4}{5.25}{3}{5.5}{1}{1}{\tfrac{1}{2}\,f_1+\tfrac{1}{4}\,g}{0}
\fxmlink{2}{4}{5.25}{4}{4.25}{2}{1}{\tfrac{1}{8}\,\bar{k}^{-}}{-1}
\fxmlink{2}{4}{5.25}{4}{4.75}{2}{1}{\tfrac{1}{8}\,\bar{k}^{-}+\tfrac{1}{2}\,k_1}{-1}
\fxmlink{2}{4}{5.25}{4}{5.75}{1}{1}{-\tfrac{1}{4}\,f_1+\tfrac{1}{4}\,f_2}{1}
\fxmlink{2}{4}{5.25}{4}{5.75}{2}{1}{\tfrac{1}{2}\,k}{-1}
\fxmlink{2}{4}{5.25}{4}{6.25}{2}{1}{\tfrac{1}{8}\,\bar{k}^{+}}{-1}
\fxmlink{2}{4}{5.25}{4}{6.75}{2}{1}{\tfrac{1}{8}\,\bar{k}^{+}}{-1}
\fxmlink{2}{4}{5.25}{4.5}{5.25}{1}{1}{\tfrac{1}{2}\,f_1+\tfrac{1}{2}\,h_1}{0}
\fxmhanging{2}{4}{4.5}
\fxmhanging{2}{4}{4.5}
\fxmhanging{2}{4}{5.5}
\fxmhanging{2}{4}{5.5}
\fxmhanging{2}{4}{6.5}
\fxmhanging{2}{4}{6.5}
\fxmcell{3}{6}{4.5}{6.5}{5}{1}{0}
\fxmcell{3}{6.5}{4.5}{7}{5}{1}{0}
\fxmcell{3}{7}{4.5}{7.5}{5}{1}{1}
\fxmsolid{3}{(7.5,4.894926) -- (7.5,5) -- (7.244029,5) -- cycle}
\fxmwall{3}{7.5}{4.894926}{7.244029}{5}
\fxmcell{3}{7.5}{4.5}{8}{5}{1}{1}
\fxmsolid{3}{(8,4.783878) -- (8,5) -- (7.5,5) -- (7.5,4.857327) -- cycle}
\fxmwall{3}{8}{4.783878}{7.5}{4.857327}
\fxmcell{3}{6}{5}{6.5}{5.5}{1}{0}
\fxmcell{3}{6.5}{5}{7}{5.5}{1}{1}
\fxmsolid{3}{(7,5.165022) -- (7,5.5) -- (6.774767,5.5) -- cycle}
\fxmwall{3}{7}{5.165022}{6.774767}{5.5}
\fxmcell{3}{7}{5}{7.5}{5.5}{1}{1}
\fxmsolid{3}{(7.22103,5) -- (7.5,5) -- (7.5,5.5) -- (7,5.5) -- (7,5.238627) -- cycle}
\fxmwall{3}{7.22103}{5}{7}{5.238627}
\fxmcell{3}{7.5}{5}{8}{5.5}{1}{2}
\fxmsolid{3}{(7.5,5) -- (8,5) -- (8,5.5) -- (7.5,5.5) -- cycle}
\fxmcell{3}{6}{5.5}{6.5}{6}{1}{0}
\fxmcell{3}{6.5}{5.5}{7}{6}{1}{1}
\fxmsolid{3}{(6.753218,5.5) -- (7,5.5) -- (7,6) -- (6.528991,6) -- cycle}
\fxmwall{3}{6.753218}{5.5}{6.528991}{6}
\fxmcell{3}{7}{5.5}{7.5}{6}{1}{2}
\fxmsolid{3}{(7,5.5) -- (7.5,5.5) -- (7.5,6) -- (7,6) -- cycle}
\fxmcell{3}{7.5}{5.5}{8}{6}{1}{2}
\fxmsolid{3}{(7.5,5.5) -- (8,5.5) -- (8,6) -- (7.5,6) -- cycle}
\fxmopen{3}{(6.5,5) -- (7,5) -- (7,5.165022) -- (6.774767,5.5) -- (6.5,5.5) -- cycle}
\fxmopen{3}{(7,5) -- (7.22103,5) -- (7,5.238627) -- cycle}
\fxmdual{3}{6.75}{5}{7.25}{5.5}{1}
\fxmface{3}{2}{7.25}{5}{0.5}{l'}{0.47}
\fxmface{3}{2}{6.75}{5.5}{0.5}{k}{0.53}
\fxmface{3}{2}{7.25}{5.5}{0.5}{k'}{0.00}
\fxmface{3}{1}{7.5}{5.25}{0.5}{g'}{0.00}
\fxmface{3}{1}{6.5}{5.25}{0.5}{g}{1.00}
\fxmface{3}{1}{7}{5.25}{0.5}{f}{0.40}
\fxmface{3}{2}{6.75}{5}{0.5}{l}{1.00}
\fxmdualface{3}{6.75}{5}{6.75}{5.5}{\tfrac12(\aperture{f}\,f+\aperture{g}\,g)}{l}
\fxmdualface{3}{6.75}{5}{7.25}{5}{\tfrac12(\aperture{l}\,l+\aperture{l'}\,l')}{b}
\fxmdualface{3}{6.75}{5.5}{7.25}{5.5}{\tfrac12(\aperture{k}\,k+\aperture{k'}\,k')}{t}
\fxmdualface{3}{7.25}{5}{7.25}{5.5}{\tfrac12(\aperture{f}\,f+\aperture{g'}\,g')}{r}
%
    \fxmnone
    \def\fxmcell##1##2##3##4##5##6##7{\ifnum##1=#1 \draw[tr cell] (##2,##3) rectangle (##4,##5);\fi}%
    \def\fxmwall##1##2##3##4##5{\ifnum##1=#1 \draw[tr solid line] (##2,##3) -- (##4,##5);\fi}%
    \def\fxmdual##1##2##3##4##5##6{\ifnum##1=#1
      \ifnum##6=1 \draw[tr dual] (##2+0.03,##3+0.03) rectangle (##4-0.03,##5-0.03);
      \else\ifnum#1=2 \draw[tr dual,line width=0.35pt,draw=black!45] (##2+0.03,##3+0.03) rectangle (##4-0.03,##5-0.03);\fi\fi\fi}%
    % Generated by docs/techref/figures/generate.jl; do not edit.
\fxmcell{1}{4}{4}{5}{5}{0}{0}
\fxmcell{1}{5}{4}{6}{5}{0}{0}
\fxmcell{1}{6}{4}{7}{5}{0}{0}
\fxmcell{1}{7}{4}{8}{5}{0}{0}
\fxmcell{1}{4}{5}{5}{6}{0}{0}
\fxmcell{1}{5}{5}{6}{6}{0}{0}
\fxmcell{1}{6}{5}{7}{6}{0}{0}
\fxmcell{1}{7}{5}{8}{6}{0}{0}
\fxmcell{1}{4}{6}{5}{7}{0}{0}
\fxmcell{1}{5}{6}{6}{7}{0}{0}
\fxmcell{1}{6}{6}{7}{7}{0}{0}
\fxmcell{1}{7}{6}{8}{7}{0}{0}
\fxmdual{1}{4.5}{4}{5.5}{5}{0}
\fxmdual{1}{4.5}{5}{5.5}{6}{0}
\fxmdual{1}{4.5}{6}{5.5}{7}{0}
\fxmdual{1}{5.5}{4}{6.5}{5}{0}
\fxmdual{1}{5.5}{5}{6.5}{6}{1}
\fxmdual{1}{5.5}{6}{6.5}{7}{0}
\fxmdual{1}{6.5}{4}{7.5}{5}{0}
\fxmdual{1}{6.5}{5}{7.5}{6}{0}
\fxmdual{1}{6.5}{6}{7.5}{7}{0}
\fxmface{1}{1}{6}{5.5}{1}{f}{1.00}
\fxmface{1}{2}{6.5}{6}{1}{k'}{1.00}
\fxmface{1}{2}{5.5}{5}{1}{l}{1.00}
\fxmface{1}{2}{6.5}{5}{1}{l'}{1.00}
\fxmface{1}{1}{7}{5.5}{1}{g'}{1.00}
\fxmface{1}{1}{5}{5.5}{1}{g}{1.00}
\fxmface{1}{2}{5.5}{6}{1}{k}{1.00}
\fxmdualface{1}{5.5}{5}{5.5}{6}{\tfrac{1}{2}\,f+\tfrac{1}{2}\,g}{l}
\fxmdualface{1}{6.5}{5}{6.5}{6}{\tfrac{1}{2}\,f+\tfrac{1}{2}\,g'}{r}
\fxmdualface{1}{5.5}{5}{6.5}{5}{\tfrac{1}{2}\,l+\tfrac{1}{2}\,l'}{b}
\fxmdualface{1}{5.5}{6}{6.5}{6}{\tfrac{1}{2}\,k+\tfrac{1}{2}\,k'}{t}
\fxmcell{2}{2}{3}{3}{4}{0}{0}
\fxmcell{2}{3}{3}{4}{4}{0}{0}
\fxmcell{2}{4}{3}{5}{4}{0}{0}
\fxmcell{2}{5}{3}{6}{4}{0}{0}
\fxmcell{2}{2}{4}{3}{5}{0}{0}
\fxmcell{2}{3}{4}{4}{5}{0}{0}
\fxmcell{2}{2}{5}{3}{6}{0}{0}
\fxmcell{2}{3}{5}{4}{6}{0}{0}
\fxmcell{2}{2}{6}{3}{7}{0}{0}
\fxmcell{2}{3}{6}{4}{7}{0}{0}
\fxmcell{2}{2}{7}{3}{8}{0}{0}
\fxmcell{2}{3}{7}{4}{8}{0}{0}
\fxmcell{2}{4}{4}{4.5}{4.5}{1}{0}
\fxmcell{2}{4.5}{4}{5}{4.5}{1}{0}
\fxmcell{2}{5}{4}{5.5}{4.5}{1}{0}
\fxmcell{2}{4}{4.5}{4.5}{5}{1}{0}
\fxmcell{2}{4.5}{4.5}{5}{5}{1}{0}
\fxmcell{2}{5}{4.5}{5.5}{5}{1}{0}
\fxmcell{2}{4}{5}{4.5}{5.5}{1}{0}
\fxmcell{2}{4.5}{5}{5}{5.5}{1}{0}
\fxmcell{2}{5}{5}{5.5}{5.5}{1}{0}
\fxmcell{2}{4}{5.5}{4.5}{6}{1}{0}
\fxmcell{2}{4.5}{5.5}{5}{6}{1}{0}
\fxmcell{2}{5}{5.5}{5.5}{6}{1}{0}
\fxmcell{2}{4}{6}{4.5}{6.5}{1}{0}
\fxmcell{2}{4.5}{6}{5}{6.5}{1}{0}
\fxmcell{2}{5}{6}{5.5}{6.5}{1}{0}
\fxmcell{2}{4}{6.5}{4.5}{7}{1}{0}
\fxmcell{2}{4.5}{6.5}{5}{7}{1}{0}
\fxmcell{2}{5}{6.5}{5.5}{7}{1}{0}
\fxmcell{2}{4}{7}{4.5}{7.5}{1}{0}
\fxmcell{2}{4.5}{7}{5}{7.5}{1}{0}
\fxmcell{2}{5}{7}{5.5}{7.5}{1}{0}
\fxmdual{2}{1.5}{3}{2.5}{4}{0}
\fxmdual{2}{1.5}{4}{2.5}{5}{0}
\fxmdual{2}{1.5}{5}{2.5}{6}{0}
\fxmdual{2}{1.5}{6}{2.5}{7}{0}
\fxmdual{2}{1.5}{7}{2.5}{8}{0}
\fxmdual{2}{2.5}{3}{3.5}{4}{0}
\fxmdual{2}{2.5}{4}{3.5}{5}{0}
\fxmdual{2}{2.5}{5}{3.5}{6}{0}
\fxmdual{2}{2.5}{6}{3.5}{7}{0}
\fxmdual{2}{2.5}{7}{3.5}{8}{0}
\fxmdual{2}{3.5}{3}{4.5}{4}{0}
\fxmdual{2}{3.5}{4}{4.25}{4.5}{0}
\fxmdual{2}{3.5}{4.5}{4.25}{5}{0}
\fxmdual{2}{3.5}{5}{4.25}{5.5}{1}
\fxmdual{2}{3.5}{5.5}{4.25}{6}{0}
\fxmdual{2}{3.5}{6}{4.25}{6.5}{0}
\fxmdual{2}{3.5}{6.5}{4.25}{7}{0}
\fxmdual{2}{3.5}{7}{4.25}{7.5}{0}
\fxmdual{2}{4.25}{4}{4.75}{4.5}{0}
\fxmdual{2}{4.25}{4.5}{4.75}{5}{0}
\fxmdual{2}{4.25}{5}{4.75}{5.5}{0}
\fxmdual{2}{4.25}{5.5}{4.75}{6}{0}
\fxmdual{2}{4.25}{6}{4.75}{6.5}{0}
\fxmdual{2}{4.25}{6.5}{4.75}{7}{0}
\fxmdual{2}{4.25}{7}{4.75}{7.5}{0}
\fxmdual{2}{4.5}{3}{5.5}{4}{0}
\fxmdual{2}{4.75}{4}{5.25}{4.5}{0}
\fxmdual{2}{4.75}{4.5}{5.25}{5}{0}
\fxmdual{2}{4.75}{5}{5.25}{5.5}{0}
\fxmdual{2}{4.75}{5.5}{5.25}{6}{0}
\fxmdual{2}{4.75}{6}{5.25}{6.5}{0}
\fxmdual{2}{4.75}{6.5}{5.25}{7}{0}
\fxmdual{2}{4.75}{7}{5.25}{7.5}{0}
\fxmface{2}{1}{4}{5.75}{0.5}{f_2}{1.00}
\fxmface{2}{1}{4}{5.25}{0.5}{f_1}{1.00}
\fxmface{2}{1}{4.5}{5.25}{0.5}{h_1}{1.00}
\fxmface{2}{2}{3.5}{6}{1}{\bar{k}^{+}}{1.00}
\fxmface{2}{2}{4.25}{5.5}{0.5}{k}{1.00}
\fxmface{2}{1}{3}{5.5}{1}{g}{1.00}
\fxmface{2}{2}{3.5}{5}{1}{\bar{k}^{-}}{1.00}
\fxmface{2}{2}{4.25}{5}{0.5}{k_1}{1.00}
\fxmlink{2}{4}{5.25}{3}{5.5}{1}{1}{\tfrac{1}{2}\,f_1+\tfrac{1}{4}\,g}{0}
\fxmlink{2}{4}{5.25}{4}{4.25}{2}{1}{\tfrac{1}{8}\,\bar{k}^{-}}{-1}
\fxmlink{2}{4}{5.25}{4}{4.75}{2}{1}{\tfrac{1}{8}\,\bar{k}^{-}+\tfrac{1}{2}\,k_1}{-1}
\fxmlink{2}{4}{5.25}{4}{5.75}{1}{1}{-\tfrac{1}{4}\,f_1+\tfrac{1}{4}\,f_2}{1}
\fxmlink{2}{4}{5.25}{4}{5.75}{2}{1}{\tfrac{1}{2}\,k}{-1}
\fxmlink{2}{4}{5.25}{4}{6.25}{2}{1}{\tfrac{1}{8}\,\bar{k}^{+}}{-1}
\fxmlink{2}{4}{5.25}{4}{6.75}{2}{1}{\tfrac{1}{8}\,\bar{k}^{+}}{-1}
\fxmlink{2}{4}{5.25}{4.5}{5.25}{1}{1}{\tfrac{1}{2}\,f_1+\tfrac{1}{2}\,h_1}{0}
\fxmhanging{2}{4}{4.5}
\fxmhanging{2}{4}{4.5}
\fxmhanging{2}{4}{5.5}
\fxmhanging{2}{4}{5.5}
\fxmhanging{2}{4}{6.5}
\fxmhanging{2}{4}{6.5}
\fxmcell{3}{6}{4.5}{6.5}{5}{1}{0}
\fxmcell{3}{6.5}{4.5}{7}{5}{1}{0}
\fxmcell{3}{7}{4.5}{7.5}{5}{1}{1}
\fxmsolid{3}{(7.5,4.894926) -- (7.5,5) -- (7.244029,5) -- cycle}
\fxmwall{3}{7.5}{4.894926}{7.244029}{5}
\fxmcell{3}{7.5}{4.5}{8}{5}{1}{1}
\fxmsolid{3}{(8,4.783878) -- (8,5) -- (7.5,5) -- (7.5,4.857327) -- cycle}
\fxmwall{3}{8}{4.783878}{7.5}{4.857327}
\fxmcell{3}{6}{5}{6.5}{5.5}{1}{0}
\fxmcell{3}{6.5}{5}{7}{5.5}{1}{1}
\fxmsolid{3}{(7,5.165022) -- (7,5.5) -- (6.774767,5.5) -- cycle}
\fxmwall{3}{7}{5.165022}{6.774767}{5.5}
\fxmcell{3}{7}{5}{7.5}{5.5}{1}{1}
\fxmsolid{3}{(7.22103,5) -- (7.5,5) -- (7.5,5.5) -- (7,5.5) -- (7,5.238627) -- cycle}
\fxmwall{3}{7.22103}{5}{7}{5.238627}
\fxmcell{3}{7.5}{5}{8}{5.5}{1}{2}
\fxmsolid{3}{(7.5,5) -- (8,5) -- (8,5.5) -- (7.5,5.5) -- cycle}
\fxmcell{3}{6}{5.5}{6.5}{6}{1}{0}
\fxmcell{3}{6.5}{5.5}{7}{6}{1}{1}
\fxmsolid{3}{(6.753218,5.5) -- (7,5.5) -- (7,6) -- (6.528991,6) -- cycle}
\fxmwall{3}{6.753218}{5.5}{6.528991}{6}
\fxmcell{3}{7}{5.5}{7.5}{6}{1}{2}
\fxmsolid{3}{(7,5.5) -- (7.5,5.5) -- (7.5,6) -- (7,6) -- cycle}
\fxmcell{3}{7.5}{5.5}{8}{6}{1}{2}
\fxmsolid{3}{(7.5,5.5) -- (8,5.5) -- (8,6) -- (7.5,6) -- cycle}
\fxmopen{3}{(6.5,5) -- (7,5) -- (7,5.165022) -- (6.774767,5.5) -- (6.5,5.5) -- cycle}
\fxmopen{3}{(7,5) -- (7.22103,5) -- (7,5.238627) -- cycle}
\fxmdual{3}{6.75}{5}{7.25}{5.5}{1}
\fxmface{3}{2}{7.25}{5}{0.5}{l'}{0.47}
\fxmface{3}{2}{6.75}{5.5}{0.5}{k}{0.53}
\fxmface{3}{2}{7.25}{5.5}{0.5}{k'}{0.00}
\fxmface{3}{1}{7.5}{5.25}{0.5}{g'}{0.00}
\fxmface{3}{1}{6.5}{5.25}{0.5}{g}{1.00}
\fxmface{3}{1}{7}{5.25}{0.5}{f}{0.40}
\fxmface{3}{2}{6.75}{5}{0.5}{l}{1.00}
\fxmdualface{3}{6.75}{5}{6.75}{5.5}{\tfrac12(\aperture{f}\,f+\aperture{g}\,g)}{l}
\fxmdualface{3}{6.75}{5}{7.25}{5}{\tfrac12(\aperture{l}\,l+\aperture{l'}\,l')}{b}
\fxmdualface{3}{6.75}{5.5}{7.25}{5.5}{\tfrac12(\aperture{k}\,k+\aperture{k'}\,k')}{t}
\fxmdualface{3}{7.25}{5}{7.25}{5.5}{\tfrac12(\aperture{f}\,f+\aperture{g'}\,g')}{r}
%
    \fxmnone
    \def\fxmface##1##2##3##4##5##6##7{\ifnum##1=#1
      \ifnum##2=1 \draw[line width=1.6pt,trred] (##3,##4-##5/2+0.04) -- (##3,##4+##5/2-0.04);
      \else \draw[line width=1.6pt,trblue] (##3-##5/2+0.04,##4) -- (##3+##5/2-0.04,##4);\fi\fi}%
    % Generated by docs/techref/figures/generate.jl; do not edit.
\fxmcell{1}{4}{4}{5}{5}{0}{0}
\fxmcell{1}{5}{4}{6}{5}{0}{0}
\fxmcell{1}{6}{4}{7}{5}{0}{0}
\fxmcell{1}{7}{4}{8}{5}{0}{0}
\fxmcell{1}{4}{5}{5}{6}{0}{0}
\fxmcell{1}{5}{5}{6}{6}{0}{0}
\fxmcell{1}{6}{5}{7}{6}{0}{0}
\fxmcell{1}{7}{5}{8}{6}{0}{0}
\fxmcell{1}{4}{6}{5}{7}{0}{0}
\fxmcell{1}{5}{6}{6}{7}{0}{0}
\fxmcell{1}{6}{6}{7}{7}{0}{0}
\fxmcell{1}{7}{6}{8}{7}{0}{0}
\fxmdual{1}{4.5}{4}{5.5}{5}{0}
\fxmdual{1}{4.5}{5}{5.5}{6}{0}
\fxmdual{1}{4.5}{6}{5.5}{7}{0}
\fxmdual{1}{5.5}{4}{6.5}{5}{0}
\fxmdual{1}{5.5}{5}{6.5}{6}{1}
\fxmdual{1}{5.5}{6}{6.5}{7}{0}
\fxmdual{1}{6.5}{4}{7.5}{5}{0}
\fxmdual{1}{6.5}{5}{7.5}{6}{0}
\fxmdual{1}{6.5}{6}{7.5}{7}{0}
\fxmface{1}{1}{6}{5.5}{1}{f}{1.00}
\fxmface{1}{2}{6.5}{6}{1}{k'}{1.00}
\fxmface{1}{2}{5.5}{5}{1}{l}{1.00}
\fxmface{1}{2}{6.5}{5}{1}{l'}{1.00}
\fxmface{1}{1}{7}{5.5}{1}{g'}{1.00}
\fxmface{1}{1}{5}{5.5}{1}{g}{1.00}
\fxmface{1}{2}{5.5}{6}{1}{k}{1.00}
\fxmdualface{1}{5.5}{5}{5.5}{6}{\tfrac{1}{2}\,f+\tfrac{1}{2}\,g}{l}
\fxmdualface{1}{6.5}{5}{6.5}{6}{\tfrac{1}{2}\,f+\tfrac{1}{2}\,g'}{r}
\fxmdualface{1}{5.5}{5}{6.5}{5}{\tfrac{1}{2}\,l+\tfrac{1}{2}\,l'}{b}
\fxmdualface{1}{5.5}{6}{6.5}{6}{\tfrac{1}{2}\,k+\tfrac{1}{2}\,k'}{t}
\fxmcell{2}{2}{3}{3}{4}{0}{0}
\fxmcell{2}{3}{3}{4}{4}{0}{0}
\fxmcell{2}{4}{3}{5}{4}{0}{0}
\fxmcell{2}{5}{3}{6}{4}{0}{0}
\fxmcell{2}{2}{4}{3}{5}{0}{0}
\fxmcell{2}{3}{4}{4}{5}{0}{0}
\fxmcell{2}{2}{5}{3}{6}{0}{0}
\fxmcell{2}{3}{5}{4}{6}{0}{0}
\fxmcell{2}{2}{6}{3}{7}{0}{0}
\fxmcell{2}{3}{6}{4}{7}{0}{0}
\fxmcell{2}{2}{7}{3}{8}{0}{0}
\fxmcell{2}{3}{7}{4}{8}{0}{0}
\fxmcell{2}{4}{4}{4.5}{4.5}{1}{0}
\fxmcell{2}{4.5}{4}{5}{4.5}{1}{0}
\fxmcell{2}{5}{4}{5.5}{4.5}{1}{0}
\fxmcell{2}{4}{4.5}{4.5}{5}{1}{0}
\fxmcell{2}{4.5}{4.5}{5}{5}{1}{0}
\fxmcell{2}{5}{4.5}{5.5}{5}{1}{0}
\fxmcell{2}{4}{5}{4.5}{5.5}{1}{0}
\fxmcell{2}{4.5}{5}{5}{5.5}{1}{0}
\fxmcell{2}{5}{5}{5.5}{5.5}{1}{0}
\fxmcell{2}{4}{5.5}{4.5}{6}{1}{0}
\fxmcell{2}{4.5}{5.5}{5}{6}{1}{0}
\fxmcell{2}{5}{5.5}{5.5}{6}{1}{0}
\fxmcell{2}{4}{6}{4.5}{6.5}{1}{0}
\fxmcell{2}{4.5}{6}{5}{6.5}{1}{0}
\fxmcell{2}{5}{6}{5.5}{6.5}{1}{0}
\fxmcell{2}{4}{6.5}{4.5}{7}{1}{0}
\fxmcell{2}{4.5}{6.5}{5}{7}{1}{0}
\fxmcell{2}{5}{6.5}{5.5}{7}{1}{0}
\fxmcell{2}{4}{7}{4.5}{7.5}{1}{0}
\fxmcell{2}{4.5}{7}{5}{7.5}{1}{0}
\fxmcell{2}{5}{7}{5.5}{7.5}{1}{0}
\fxmdual{2}{1.5}{3}{2.5}{4}{0}
\fxmdual{2}{1.5}{4}{2.5}{5}{0}
\fxmdual{2}{1.5}{5}{2.5}{6}{0}
\fxmdual{2}{1.5}{6}{2.5}{7}{0}
\fxmdual{2}{1.5}{7}{2.5}{8}{0}
\fxmdual{2}{2.5}{3}{3.5}{4}{0}
\fxmdual{2}{2.5}{4}{3.5}{5}{0}
\fxmdual{2}{2.5}{5}{3.5}{6}{0}
\fxmdual{2}{2.5}{6}{3.5}{7}{0}
\fxmdual{2}{2.5}{7}{3.5}{8}{0}
\fxmdual{2}{3.5}{3}{4.5}{4}{0}
\fxmdual{2}{3.5}{4}{4.25}{4.5}{0}
\fxmdual{2}{3.5}{4.5}{4.25}{5}{0}
\fxmdual{2}{3.5}{5}{4.25}{5.5}{1}
\fxmdual{2}{3.5}{5.5}{4.25}{6}{0}
\fxmdual{2}{3.5}{6}{4.25}{6.5}{0}
\fxmdual{2}{3.5}{6.5}{4.25}{7}{0}
\fxmdual{2}{3.5}{7}{4.25}{7.5}{0}
\fxmdual{2}{4.25}{4}{4.75}{4.5}{0}
\fxmdual{2}{4.25}{4.5}{4.75}{5}{0}
\fxmdual{2}{4.25}{5}{4.75}{5.5}{0}
\fxmdual{2}{4.25}{5.5}{4.75}{6}{0}
\fxmdual{2}{4.25}{6}{4.75}{6.5}{0}
\fxmdual{2}{4.25}{6.5}{4.75}{7}{0}
\fxmdual{2}{4.25}{7}{4.75}{7.5}{0}
\fxmdual{2}{4.5}{3}{5.5}{4}{0}
\fxmdual{2}{4.75}{4}{5.25}{4.5}{0}
\fxmdual{2}{4.75}{4.5}{5.25}{5}{0}
\fxmdual{2}{4.75}{5}{5.25}{5.5}{0}
\fxmdual{2}{4.75}{5.5}{5.25}{6}{0}
\fxmdual{2}{4.75}{6}{5.25}{6.5}{0}
\fxmdual{2}{4.75}{6.5}{5.25}{7}{0}
\fxmdual{2}{4.75}{7}{5.25}{7.5}{0}
\fxmface{2}{1}{4}{5.75}{0.5}{f_2}{1.00}
\fxmface{2}{1}{4}{5.25}{0.5}{f_1}{1.00}
\fxmface{2}{1}{4.5}{5.25}{0.5}{h_1}{1.00}
\fxmface{2}{2}{3.5}{6}{1}{\bar{k}^{+}}{1.00}
\fxmface{2}{2}{4.25}{5.5}{0.5}{k}{1.00}
\fxmface{2}{1}{3}{5.5}{1}{g}{1.00}
\fxmface{2}{2}{3.5}{5}{1}{\bar{k}^{-}}{1.00}
\fxmface{2}{2}{4.25}{5}{0.5}{k_1}{1.00}
\fxmlink{2}{4}{5.25}{3}{5.5}{1}{1}{\tfrac{1}{2}\,f_1+\tfrac{1}{4}\,g}{0}
\fxmlink{2}{4}{5.25}{4}{4.25}{2}{1}{\tfrac{1}{8}\,\bar{k}^{-}}{-1}
\fxmlink{2}{4}{5.25}{4}{4.75}{2}{1}{\tfrac{1}{8}\,\bar{k}^{-}+\tfrac{1}{2}\,k_1}{-1}
\fxmlink{2}{4}{5.25}{4}{5.75}{1}{1}{-\tfrac{1}{4}\,f_1+\tfrac{1}{4}\,f_2}{1}
\fxmlink{2}{4}{5.25}{4}{5.75}{2}{1}{\tfrac{1}{2}\,k}{-1}
\fxmlink{2}{4}{5.25}{4}{6.25}{2}{1}{\tfrac{1}{8}\,\bar{k}^{+}}{-1}
\fxmlink{2}{4}{5.25}{4}{6.75}{2}{1}{\tfrac{1}{8}\,\bar{k}^{+}}{-1}
\fxmlink{2}{4}{5.25}{4.5}{5.25}{1}{1}{\tfrac{1}{2}\,f_1+\tfrac{1}{2}\,h_1}{0}
\fxmhanging{2}{4}{4.5}
\fxmhanging{2}{4}{4.5}
\fxmhanging{2}{4}{5.5}
\fxmhanging{2}{4}{5.5}
\fxmhanging{2}{4}{6.5}
\fxmhanging{2}{4}{6.5}
\fxmcell{3}{6}{4.5}{6.5}{5}{1}{0}
\fxmcell{3}{6.5}{4.5}{7}{5}{1}{0}
\fxmcell{3}{7}{4.5}{7.5}{5}{1}{1}
\fxmsolid{3}{(7.5,4.894926) -- (7.5,5) -- (7.244029,5) -- cycle}
\fxmwall{3}{7.5}{4.894926}{7.244029}{5}
\fxmcell{3}{7.5}{4.5}{8}{5}{1}{1}
\fxmsolid{3}{(8,4.783878) -- (8,5) -- (7.5,5) -- (7.5,4.857327) -- cycle}
\fxmwall{3}{8}{4.783878}{7.5}{4.857327}
\fxmcell{3}{6}{5}{6.5}{5.5}{1}{0}
\fxmcell{3}{6.5}{5}{7}{5.5}{1}{1}
\fxmsolid{3}{(7,5.165022) -- (7,5.5) -- (6.774767,5.5) -- cycle}
\fxmwall{3}{7}{5.165022}{6.774767}{5.5}
\fxmcell{3}{7}{5}{7.5}{5.5}{1}{1}
\fxmsolid{3}{(7.22103,5) -- (7.5,5) -- (7.5,5.5) -- (7,5.5) -- (7,5.238627) -- cycle}
\fxmwall{3}{7.22103}{5}{7}{5.238627}
\fxmcell{3}{7.5}{5}{8}{5.5}{1}{2}
\fxmsolid{3}{(7.5,5) -- (8,5) -- (8,5.5) -- (7.5,5.5) -- cycle}
\fxmcell{3}{6}{5.5}{6.5}{6}{1}{0}
\fxmcell{3}{6.5}{5.5}{7}{6}{1}{1}
\fxmsolid{3}{(6.753218,5.5) -- (7,5.5) -- (7,6) -- (6.528991,6) -- cycle}
\fxmwall{3}{6.753218}{5.5}{6.528991}{6}
\fxmcell{3}{7}{5.5}{7.5}{6}{1}{2}
\fxmsolid{3}{(7,5.5) -- (7.5,5.5) -- (7.5,6) -- (7,6) -- cycle}
\fxmcell{3}{7.5}{5.5}{8}{6}{1}{2}
\fxmsolid{3}{(7.5,5.5) -- (8,5.5) -- (8,6) -- (7.5,6) -- cycle}
\fxmopen{3}{(6.5,5) -- (7,5) -- (7,5.165022) -- (6.774767,5.5) -- (6.5,5.5) -- cycle}
\fxmopen{3}{(7,5) -- (7.22103,5) -- (7,5.238627) -- cycle}
\fxmdual{3}{6.75}{5}{7.25}{5.5}{1}
\fxmface{3}{2}{7.25}{5}{0.5}{l'}{0.47}
\fxmface{3}{2}{6.75}{5.5}{0.5}{k}{0.53}
\fxmface{3}{2}{7.25}{5.5}{0.5}{k'}{0.00}
\fxmface{3}{1}{7.5}{5.25}{0.5}{g'}{0.00}
\fxmface{3}{1}{6.5}{5.25}{0.5}{g}{1.00}
\fxmface{3}{1}{7}{5.25}{0.5}{f}{0.40}
\fxmface{3}{2}{6.75}{5}{0.5}{l}{1.00}
\fxmdualface{3}{6.75}{5}{6.75}{5.5}{\tfrac12(\aperture{f}\,f+\aperture{g}\,g)}{l}
\fxmdualface{3}{6.75}{5}{7.25}{5}{\tfrac12(\aperture{l}\,l+\aperture{l'}\,l')}{b}
\fxmdualface{3}{6.75}{5.5}{7.25}{5.5}{\tfrac12(\aperture{k}\,k+\aperture{k'}\,k')}{t}
\fxmdualface{3}{7.25}{5}{7.25}{5.5}{\tfrac12(\aperture{f}\,f+\aperture{g'}\,g')}{r}
}%
  % Labels of the primal faces of panel #1: x offset #2 (negative: left) and
  % height #5 (fraction of the face) of the labels of faces normal to e_1,
  % y offset #3 of the others; #4 = 1 adds the aperture. A panel may override
  % the offset and height of one face through \fxmlabeladjust{name}, which
  % can redefine \fxmdx and \fxmfrac.
  \def\fxmlabeladjust#1{}%
  \def\fxmfacelabels#1#2#3#4#5{%
    \fxmnone
    \def\fxmface##1##2##3##4##5##6##7{\ifnum##1=#1
      \ifnum##2=1 \def\fxmdx{#2}\def\fxmfrac{#5}\fxmlabeladjust{##6}%
        \ifdim\fxmdx pt<0pt
          \node[tr small,text=trred!80!black,anchor=east,inner sep=0.6pt,fill=white,fill opacity=0.8,text opacity=1]
            at (##3+\fxmdx,##4+##5*\fxmfrac) {$##6\ifnum#4=1 {:}\,##7\fi$};
        \else
          \node[tr small,text=trred!80!black,anchor=west,inner sep=0.6pt,fill=white,fill opacity=0.8,text opacity=1]
            at (##3+\fxmdx,##4+##5*\fxmfrac) {$##6\ifnum#4=1 {:}\,##7\fi$};
        \fi
      \else \def\fxmfrac{0.3}\fxmlabeladjust{##6}%
        \node[tr small,text=trblue!80!black,anchor=north,inner sep=0.6pt,fill=white,fill opacity=0.8,text opacity=1]
          at (##3+##5*\fxmfrac,##4+#3) {$##6\ifnum#4=1 {:}\,##7\fi$};\fi\fi}%
    % Generated by docs/techref/figures/generate.jl; do not edit.
\fxmcell{1}{4}{4}{5}{5}{0}{0}
\fxmcell{1}{5}{4}{6}{5}{0}{0}
\fxmcell{1}{6}{4}{7}{5}{0}{0}
\fxmcell{1}{7}{4}{8}{5}{0}{0}
\fxmcell{1}{4}{5}{5}{6}{0}{0}
\fxmcell{1}{5}{5}{6}{6}{0}{0}
\fxmcell{1}{6}{5}{7}{6}{0}{0}
\fxmcell{1}{7}{5}{8}{6}{0}{0}
\fxmcell{1}{4}{6}{5}{7}{0}{0}
\fxmcell{1}{5}{6}{6}{7}{0}{0}
\fxmcell{1}{6}{6}{7}{7}{0}{0}
\fxmcell{1}{7}{6}{8}{7}{0}{0}
\fxmdual{1}{4.5}{4}{5.5}{5}{0}
\fxmdual{1}{4.5}{5}{5.5}{6}{0}
\fxmdual{1}{4.5}{6}{5.5}{7}{0}
\fxmdual{1}{5.5}{4}{6.5}{5}{0}
\fxmdual{1}{5.5}{5}{6.5}{6}{1}
\fxmdual{1}{5.5}{6}{6.5}{7}{0}
\fxmdual{1}{6.5}{4}{7.5}{5}{0}
\fxmdual{1}{6.5}{5}{7.5}{6}{0}
\fxmdual{1}{6.5}{6}{7.5}{7}{0}
\fxmface{1}{1}{6}{5.5}{1}{f}{1.00}
\fxmface{1}{2}{6.5}{6}{1}{k'}{1.00}
\fxmface{1}{2}{5.5}{5}{1}{l}{1.00}
\fxmface{1}{2}{6.5}{5}{1}{l'}{1.00}
\fxmface{1}{1}{7}{5.5}{1}{g'}{1.00}
\fxmface{1}{1}{5}{5.5}{1}{g}{1.00}
\fxmface{1}{2}{5.5}{6}{1}{k}{1.00}
\fxmdualface{1}{5.5}{5}{5.5}{6}{\tfrac{1}{2}\,f+\tfrac{1}{2}\,g}{l}
\fxmdualface{1}{6.5}{5}{6.5}{6}{\tfrac{1}{2}\,f+\tfrac{1}{2}\,g'}{r}
\fxmdualface{1}{5.5}{5}{6.5}{5}{\tfrac{1}{2}\,l+\tfrac{1}{2}\,l'}{b}
\fxmdualface{1}{5.5}{6}{6.5}{6}{\tfrac{1}{2}\,k+\tfrac{1}{2}\,k'}{t}
\fxmcell{2}{2}{3}{3}{4}{0}{0}
\fxmcell{2}{3}{3}{4}{4}{0}{0}
\fxmcell{2}{4}{3}{5}{4}{0}{0}
\fxmcell{2}{5}{3}{6}{4}{0}{0}
\fxmcell{2}{2}{4}{3}{5}{0}{0}
\fxmcell{2}{3}{4}{4}{5}{0}{0}
\fxmcell{2}{2}{5}{3}{6}{0}{0}
\fxmcell{2}{3}{5}{4}{6}{0}{0}
\fxmcell{2}{2}{6}{3}{7}{0}{0}
\fxmcell{2}{3}{6}{4}{7}{0}{0}
\fxmcell{2}{2}{7}{3}{8}{0}{0}
\fxmcell{2}{3}{7}{4}{8}{0}{0}
\fxmcell{2}{4}{4}{4.5}{4.5}{1}{0}
\fxmcell{2}{4.5}{4}{5}{4.5}{1}{0}
\fxmcell{2}{5}{4}{5.5}{4.5}{1}{0}
\fxmcell{2}{4}{4.5}{4.5}{5}{1}{0}
\fxmcell{2}{4.5}{4.5}{5}{5}{1}{0}
\fxmcell{2}{5}{4.5}{5.5}{5}{1}{0}
\fxmcell{2}{4}{5}{4.5}{5.5}{1}{0}
\fxmcell{2}{4.5}{5}{5}{5.5}{1}{0}
\fxmcell{2}{5}{5}{5.5}{5.5}{1}{0}
\fxmcell{2}{4}{5.5}{4.5}{6}{1}{0}
\fxmcell{2}{4.5}{5.5}{5}{6}{1}{0}
\fxmcell{2}{5}{5.5}{5.5}{6}{1}{0}
\fxmcell{2}{4}{6}{4.5}{6.5}{1}{0}
\fxmcell{2}{4.5}{6}{5}{6.5}{1}{0}
\fxmcell{2}{5}{6}{5.5}{6.5}{1}{0}
\fxmcell{2}{4}{6.5}{4.5}{7}{1}{0}
\fxmcell{2}{4.5}{6.5}{5}{7}{1}{0}
\fxmcell{2}{5}{6.5}{5.5}{7}{1}{0}
\fxmcell{2}{4}{7}{4.5}{7.5}{1}{0}
\fxmcell{2}{4.5}{7}{5}{7.5}{1}{0}
\fxmcell{2}{5}{7}{5.5}{7.5}{1}{0}
\fxmdual{2}{1.5}{3}{2.5}{4}{0}
\fxmdual{2}{1.5}{4}{2.5}{5}{0}
\fxmdual{2}{1.5}{5}{2.5}{6}{0}
\fxmdual{2}{1.5}{6}{2.5}{7}{0}
\fxmdual{2}{1.5}{7}{2.5}{8}{0}
\fxmdual{2}{2.5}{3}{3.5}{4}{0}
\fxmdual{2}{2.5}{4}{3.5}{5}{0}
\fxmdual{2}{2.5}{5}{3.5}{6}{0}
\fxmdual{2}{2.5}{6}{3.5}{7}{0}
\fxmdual{2}{2.5}{7}{3.5}{8}{0}
\fxmdual{2}{3.5}{3}{4.5}{4}{0}
\fxmdual{2}{3.5}{4}{4.25}{4.5}{0}
\fxmdual{2}{3.5}{4.5}{4.25}{5}{0}
\fxmdual{2}{3.5}{5}{4.25}{5.5}{1}
\fxmdual{2}{3.5}{5.5}{4.25}{6}{0}
\fxmdual{2}{3.5}{6}{4.25}{6.5}{0}
\fxmdual{2}{3.5}{6.5}{4.25}{7}{0}
\fxmdual{2}{3.5}{7}{4.25}{7.5}{0}
\fxmdual{2}{4.25}{4}{4.75}{4.5}{0}
\fxmdual{2}{4.25}{4.5}{4.75}{5}{0}
\fxmdual{2}{4.25}{5}{4.75}{5.5}{0}
\fxmdual{2}{4.25}{5.5}{4.75}{6}{0}
\fxmdual{2}{4.25}{6}{4.75}{6.5}{0}
\fxmdual{2}{4.25}{6.5}{4.75}{7}{0}
\fxmdual{2}{4.25}{7}{4.75}{7.5}{0}
\fxmdual{2}{4.5}{3}{5.5}{4}{0}
\fxmdual{2}{4.75}{4}{5.25}{4.5}{0}
\fxmdual{2}{4.75}{4.5}{5.25}{5}{0}
\fxmdual{2}{4.75}{5}{5.25}{5.5}{0}
\fxmdual{2}{4.75}{5.5}{5.25}{6}{0}
\fxmdual{2}{4.75}{6}{5.25}{6.5}{0}
\fxmdual{2}{4.75}{6.5}{5.25}{7}{0}
\fxmdual{2}{4.75}{7}{5.25}{7.5}{0}
\fxmface{2}{1}{4}{5.75}{0.5}{f_2}{1.00}
\fxmface{2}{1}{4}{5.25}{0.5}{f_1}{1.00}
\fxmface{2}{1}{4.5}{5.25}{0.5}{h_1}{1.00}
\fxmface{2}{2}{3.5}{6}{1}{\bar{k}^{+}}{1.00}
\fxmface{2}{2}{4.25}{5.5}{0.5}{k}{1.00}
\fxmface{2}{1}{3}{5.5}{1}{g}{1.00}
\fxmface{2}{2}{3.5}{5}{1}{\bar{k}^{-}}{1.00}
\fxmface{2}{2}{4.25}{5}{0.5}{k_1}{1.00}
\fxmlink{2}{4}{5.25}{3}{5.5}{1}{1}{\tfrac{1}{2}\,f_1+\tfrac{1}{4}\,g}{0}
\fxmlink{2}{4}{5.25}{4}{4.25}{2}{1}{\tfrac{1}{8}\,\bar{k}^{-}}{-1}
\fxmlink{2}{4}{5.25}{4}{4.75}{2}{1}{\tfrac{1}{8}\,\bar{k}^{-}+\tfrac{1}{2}\,k_1}{-1}
\fxmlink{2}{4}{5.25}{4}{5.75}{1}{1}{-\tfrac{1}{4}\,f_1+\tfrac{1}{4}\,f_2}{1}
\fxmlink{2}{4}{5.25}{4}{5.75}{2}{1}{\tfrac{1}{2}\,k}{-1}
\fxmlink{2}{4}{5.25}{4}{6.25}{2}{1}{\tfrac{1}{8}\,\bar{k}^{+}}{-1}
\fxmlink{2}{4}{5.25}{4}{6.75}{2}{1}{\tfrac{1}{8}\,\bar{k}^{+}}{-1}
\fxmlink{2}{4}{5.25}{4.5}{5.25}{1}{1}{\tfrac{1}{2}\,f_1+\tfrac{1}{2}\,h_1}{0}
\fxmhanging{2}{4}{4.5}
\fxmhanging{2}{4}{4.5}
\fxmhanging{2}{4}{5.5}
\fxmhanging{2}{4}{5.5}
\fxmhanging{2}{4}{6.5}
\fxmhanging{2}{4}{6.5}
\fxmcell{3}{6}{4.5}{6.5}{5}{1}{0}
\fxmcell{3}{6.5}{4.5}{7}{5}{1}{0}
\fxmcell{3}{7}{4.5}{7.5}{5}{1}{1}
\fxmsolid{3}{(7.5,4.894926) -- (7.5,5) -- (7.244029,5) -- cycle}
\fxmwall{3}{7.5}{4.894926}{7.244029}{5}
\fxmcell{3}{7.5}{4.5}{8}{5}{1}{1}
\fxmsolid{3}{(8,4.783878) -- (8,5) -- (7.5,5) -- (7.5,4.857327) -- cycle}
\fxmwall{3}{8}{4.783878}{7.5}{4.857327}
\fxmcell{3}{6}{5}{6.5}{5.5}{1}{0}
\fxmcell{3}{6.5}{5}{7}{5.5}{1}{1}
\fxmsolid{3}{(7,5.165022) -- (7,5.5) -- (6.774767,5.5) -- cycle}
\fxmwall{3}{7}{5.165022}{6.774767}{5.5}
\fxmcell{3}{7}{5}{7.5}{5.5}{1}{1}
\fxmsolid{3}{(7.22103,5) -- (7.5,5) -- (7.5,5.5) -- (7,5.5) -- (7,5.238627) -- cycle}
\fxmwall{3}{7.22103}{5}{7}{5.238627}
\fxmcell{3}{7.5}{5}{8}{5.5}{1}{2}
\fxmsolid{3}{(7.5,5) -- (8,5) -- (8,5.5) -- (7.5,5.5) -- cycle}
\fxmcell{3}{6}{5.5}{6.5}{6}{1}{0}
\fxmcell{3}{6.5}{5.5}{7}{6}{1}{1}
\fxmsolid{3}{(6.753218,5.5) -- (7,5.5) -- (7,6) -- (6.528991,6) -- cycle}
\fxmwall{3}{6.753218}{5.5}{6.528991}{6}
\fxmcell{3}{7}{5.5}{7.5}{6}{1}{2}
\fxmsolid{3}{(7,5.5) -- (7.5,5.5) -- (7.5,6) -- (7,6) -- cycle}
\fxmcell{3}{7.5}{5.5}{8}{6}{1}{2}
\fxmsolid{3}{(7.5,5.5) -- (8,5.5) -- (8,6) -- (7.5,6) -- cycle}
\fxmopen{3}{(6.5,5) -- (7,5) -- (7,5.165022) -- (6.774767,5.5) -- (6.5,5.5) -- cycle}
\fxmopen{3}{(7,5) -- (7.22103,5) -- (7,5.238627) -- cycle}
\fxmdual{3}{6.75}{5}{7.25}{5.5}{1}
\fxmface{3}{2}{7.25}{5}{0.5}{l'}{0.47}
\fxmface{3}{2}{6.75}{5.5}{0.5}{k}{0.53}
\fxmface{3}{2}{7.25}{5.5}{0.5}{k'}{0.00}
\fxmface{3}{1}{7.5}{5.25}{0.5}{g'}{0.00}
\fxmface{3}{1}{6.5}{5.25}{0.5}{g}{1.00}
\fxmface{3}{1}{7}{5.25}{0.5}{f}{0.40}
\fxmface{3}{2}{6.75}{5}{0.5}{l}{1.00}
\fxmdualface{3}{6.75}{5}{6.75}{5.5}{\tfrac12(\aperture{f}\,f+\aperture{g}\,g)}{l}
\fxmdualface{3}{6.75}{5}{7.25}{5}{\tfrac12(\aperture{l}\,l+\aperture{l'}\,l')}{b}
\fxmdualface{3}{6.75}{5.5}{7.25}{5.5}{\tfrac12(\aperture{k}\,k+\aperture{k'}\,k')}{t}
\fxmdualface{3}{7.25}{5}{7.25}{5.5}{\tfrac12(\aperture{f}\,f+\aperture{g'}\,g')}{r}
}%
  \begin{minipage}[b]{0.46\textwidth}
  \begin{subfigure}{\linewidth}
    \centering
    \begin{tikzpicture}[x=1.95cm,y=1.95cm]
      \begin{scope}
        \clip (5.05,4.6) rectangle (6.95,6.4);
        \fxmpasses{1}
      \end{scope}
      \fxmfacelabels{1}{0.03}{-0.03}{0}{0.3}
      \fxmnone
      \def\fxmdualface#1#2#3#4#5#6#7{\ifnum#1=1
        \draw[trpurple,line width=1.1pt] (#2,#3) -- (#4,#5);
        \if l#7\node[tr link label,anchor=east] at (#2-0.05,{(#3+#5)/2-0.25}) {$#6$};\fi
        \if r#7\node[tr link label,anchor=west] at (#2+0.05,{(#3+#5)/2-0.25}) {$#6$};\fi
        \if b#7\node[tr link label,anchor=north] at ({(#2+#4)/2},#3-0.12) {$#6$};\fi
        \if t#7\node[tr link label,anchor=south] at ({(#2+#4)/2},#3+0.05) {$#6$};\fi\fi}%
      % Generated by docs/techref/figures/generate.jl; do not edit.
\fxmcell{1}{4}{4}{5}{5}{0}{0}
\fxmcell{1}{5}{4}{6}{5}{0}{0}
\fxmcell{1}{6}{4}{7}{5}{0}{0}
\fxmcell{1}{7}{4}{8}{5}{0}{0}
\fxmcell{1}{4}{5}{5}{6}{0}{0}
\fxmcell{1}{5}{5}{6}{6}{0}{0}
\fxmcell{1}{6}{5}{7}{6}{0}{0}
\fxmcell{1}{7}{5}{8}{6}{0}{0}
\fxmcell{1}{4}{6}{5}{7}{0}{0}
\fxmcell{1}{5}{6}{6}{7}{0}{0}
\fxmcell{1}{6}{6}{7}{7}{0}{0}
\fxmcell{1}{7}{6}{8}{7}{0}{0}
\fxmdual{1}{4.5}{4}{5.5}{5}{0}
\fxmdual{1}{4.5}{5}{5.5}{6}{0}
\fxmdual{1}{4.5}{6}{5.5}{7}{0}
\fxmdual{1}{5.5}{4}{6.5}{5}{0}
\fxmdual{1}{5.5}{5}{6.5}{6}{1}
\fxmdual{1}{5.5}{6}{6.5}{7}{0}
\fxmdual{1}{6.5}{4}{7.5}{5}{0}
\fxmdual{1}{6.5}{5}{7.5}{6}{0}
\fxmdual{1}{6.5}{6}{7.5}{7}{0}
\fxmface{1}{1}{6}{5.5}{1}{f}{1.00}
\fxmface{1}{2}{6.5}{6}{1}{k'}{1.00}
\fxmface{1}{2}{5.5}{5}{1}{l}{1.00}
\fxmface{1}{2}{6.5}{5}{1}{l'}{1.00}
\fxmface{1}{1}{7}{5.5}{1}{g'}{1.00}
\fxmface{1}{1}{5}{5.5}{1}{g}{1.00}
\fxmface{1}{2}{5.5}{6}{1}{k}{1.00}
\fxmdualface{1}{5.5}{5}{5.5}{6}{\tfrac{1}{2}\,f+\tfrac{1}{2}\,g}{l}
\fxmdualface{1}{6.5}{5}{6.5}{6}{\tfrac{1}{2}\,f+\tfrac{1}{2}\,g'}{r}
\fxmdualface{1}{5.5}{5}{6.5}{5}{\tfrac{1}{2}\,l+\tfrac{1}{2}\,l'}{b}
\fxmdualface{1}{5.5}{6}{6.5}{6}{\tfrac{1}{2}\,k+\tfrac{1}{2}\,k'}{t}
\fxmcell{2}{2}{3}{3}{4}{0}{0}
\fxmcell{2}{3}{3}{4}{4}{0}{0}
\fxmcell{2}{4}{3}{5}{4}{0}{0}
\fxmcell{2}{5}{3}{6}{4}{0}{0}
\fxmcell{2}{2}{4}{3}{5}{0}{0}
\fxmcell{2}{3}{4}{4}{5}{0}{0}
\fxmcell{2}{2}{5}{3}{6}{0}{0}
\fxmcell{2}{3}{5}{4}{6}{0}{0}
\fxmcell{2}{2}{6}{3}{7}{0}{0}
\fxmcell{2}{3}{6}{4}{7}{0}{0}
\fxmcell{2}{2}{7}{3}{8}{0}{0}
\fxmcell{2}{3}{7}{4}{8}{0}{0}
\fxmcell{2}{4}{4}{4.5}{4.5}{1}{0}
\fxmcell{2}{4.5}{4}{5}{4.5}{1}{0}
\fxmcell{2}{5}{4}{5.5}{4.5}{1}{0}
\fxmcell{2}{4}{4.5}{4.5}{5}{1}{0}
\fxmcell{2}{4.5}{4.5}{5}{5}{1}{0}
\fxmcell{2}{5}{4.5}{5.5}{5}{1}{0}
\fxmcell{2}{4}{5}{4.5}{5.5}{1}{0}
\fxmcell{2}{4.5}{5}{5}{5.5}{1}{0}
\fxmcell{2}{5}{5}{5.5}{5.5}{1}{0}
\fxmcell{2}{4}{5.5}{4.5}{6}{1}{0}
\fxmcell{2}{4.5}{5.5}{5}{6}{1}{0}
\fxmcell{2}{5}{5.5}{5.5}{6}{1}{0}
\fxmcell{2}{4}{6}{4.5}{6.5}{1}{0}
\fxmcell{2}{4.5}{6}{5}{6.5}{1}{0}
\fxmcell{2}{5}{6}{5.5}{6.5}{1}{0}
\fxmcell{2}{4}{6.5}{4.5}{7}{1}{0}
\fxmcell{2}{4.5}{6.5}{5}{7}{1}{0}
\fxmcell{2}{5}{6.5}{5.5}{7}{1}{0}
\fxmcell{2}{4}{7}{4.5}{7.5}{1}{0}
\fxmcell{2}{4.5}{7}{5}{7.5}{1}{0}
\fxmcell{2}{5}{7}{5.5}{7.5}{1}{0}
\fxmdual{2}{1.5}{3}{2.5}{4}{0}
\fxmdual{2}{1.5}{4}{2.5}{5}{0}
\fxmdual{2}{1.5}{5}{2.5}{6}{0}
\fxmdual{2}{1.5}{6}{2.5}{7}{0}
\fxmdual{2}{1.5}{7}{2.5}{8}{0}
\fxmdual{2}{2.5}{3}{3.5}{4}{0}
\fxmdual{2}{2.5}{4}{3.5}{5}{0}
\fxmdual{2}{2.5}{5}{3.5}{6}{0}
\fxmdual{2}{2.5}{6}{3.5}{7}{0}
\fxmdual{2}{2.5}{7}{3.5}{8}{0}
\fxmdual{2}{3.5}{3}{4.5}{4}{0}
\fxmdual{2}{3.5}{4}{4.25}{4.5}{0}
\fxmdual{2}{3.5}{4.5}{4.25}{5}{0}
\fxmdual{2}{3.5}{5}{4.25}{5.5}{1}
\fxmdual{2}{3.5}{5.5}{4.25}{6}{0}
\fxmdual{2}{3.5}{6}{4.25}{6.5}{0}
\fxmdual{2}{3.5}{6.5}{4.25}{7}{0}
\fxmdual{2}{3.5}{7}{4.25}{7.5}{0}
\fxmdual{2}{4.25}{4}{4.75}{4.5}{0}
\fxmdual{2}{4.25}{4.5}{4.75}{5}{0}
\fxmdual{2}{4.25}{5}{4.75}{5.5}{0}
\fxmdual{2}{4.25}{5.5}{4.75}{6}{0}
\fxmdual{2}{4.25}{6}{4.75}{6.5}{0}
\fxmdual{2}{4.25}{6.5}{4.75}{7}{0}
\fxmdual{2}{4.25}{7}{4.75}{7.5}{0}
\fxmdual{2}{4.5}{3}{5.5}{4}{0}
\fxmdual{2}{4.75}{4}{5.25}{4.5}{0}
\fxmdual{2}{4.75}{4.5}{5.25}{5}{0}
\fxmdual{2}{4.75}{5}{5.25}{5.5}{0}
\fxmdual{2}{4.75}{5.5}{5.25}{6}{0}
\fxmdual{2}{4.75}{6}{5.25}{6.5}{0}
\fxmdual{2}{4.75}{6.5}{5.25}{7}{0}
\fxmdual{2}{4.75}{7}{5.25}{7.5}{0}
\fxmface{2}{1}{4}{5.75}{0.5}{f_2}{1.00}
\fxmface{2}{1}{4}{5.25}{0.5}{f_1}{1.00}
\fxmface{2}{1}{4.5}{5.25}{0.5}{h_1}{1.00}
\fxmface{2}{2}{3.5}{6}{1}{\bar{k}^{+}}{1.00}
\fxmface{2}{2}{4.25}{5.5}{0.5}{k}{1.00}
\fxmface{2}{1}{3}{5.5}{1}{g}{1.00}
\fxmface{2}{2}{3.5}{5}{1}{\bar{k}^{-}}{1.00}
\fxmface{2}{2}{4.25}{5}{0.5}{k_1}{1.00}
\fxmlink{2}{4}{5.25}{3}{5.5}{1}{1}{\tfrac{1}{2}\,f_1+\tfrac{1}{4}\,g}{0}
\fxmlink{2}{4}{5.25}{4}{4.25}{2}{1}{\tfrac{1}{8}\,\bar{k}^{-}}{-1}
\fxmlink{2}{4}{5.25}{4}{4.75}{2}{1}{\tfrac{1}{8}\,\bar{k}^{-}+\tfrac{1}{2}\,k_1}{-1}
\fxmlink{2}{4}{5.25}{4}{5.75}{1}{1}{-\tfrac{1}{4}\,f_1+\tfrac{1}{4}\,f_2}{1}
\fxmlink{2}{4}{5.25}{4}{5.75}{2}{1}{\tfrac{1}{2}\,k}{-1}
\fxmlink{2}{4}{5.25}{4}{6.25}{2}{1}{\tfrac{1}{8}\,\bar{k}^{+}}{-1}
\fxmlink{2}{4}{5.25}{4}{6.75}{2}{1}{\tfrac{1}{8}\,\bar{k}^{+}}{-1}
\fxmlink{2}{4}{5.25}{4.5}{5.25}{1}{1}{\tfrac{1}{2}\,f_1+\tfrac{1}{2}\,h_1}{0}
\fxmhanging{2}{4}{4.5}
\fxmhanging{2}{4}{4.5}
\fxmhanging{2}{4}{5.5}
\fxmhanging{2}{4}{5.5}
\fxmhanging{2}{4}{6.5}
\fxmhanging{2}{4}{6.5}
\fxmcell{3}{6}{4.5}{6.5}{5}{1}{0}
\fxmcell{3}{6.5}{4.5}{7}{5}{1}{0}
\fxmcell{3}{7}{4.5}{7.5}{5}{1}{1}
\fxmsolid{3}{(7.5,4.894926) -- (7.5,5) -- (7.244029,5) -- cycle}
\fxmwall{3}{7.5}{4.894926}{7.244029}{5}
\fxmcell{3}{7.5}{4.5}{8}{5}{1}{1}
\fxmsolid{3}{(8,4.783878) -- (8,5) -- (7.5,5) -- (7.5,4.857327) -- cycle}
\fxmwall{3}{8}{4.783878}{7.5}{4.857327}
\fxmcell{3}{6}{5}{6.5}{5.5}{1}{0}
\fxmcell{3}{6.5}{5}{7}{5.5}{1}{1}
\fxmsolid{3}{(7,5.165022) -- (7,5.5) -- (6.774767,5.5) -- cycle}
\fxmwall{3}{7}{5.165022}{6.774767}{5.5}
\fxmcell{3}{7}{5}{7.5}{5.5}{1}{1}
\fxmsolid{3}{(7.22103,5) -- (7.5,5) -- (7.5,5.5) -- (7,5.5) -- (7,5.238627) -- cycle}
\fxmwall{3}{7.22103}{5}{7}{5.238627}
\fxmcell{3}{7.5}{5}{8}{5.5}{1}{2}
\fxmsolid{3}{(7.5,5) -- (8,5) -- (8,5.5) -- (7.5,5.5) -- cycle}
\fxmcell{3}{6}{5.5}{6.5}{6}{1}{0}
\fxmcell{3}{6.5}{5.5}{7}{6}{1}{1}
\fxmsolid{3}{(6.753218,5.5) -- (7,5.5) -- (7,6) -- (6.528991,6) -- cycle}
\fxmwall{3}{6.753218}{5.5}{6.528991}{6}
\fxmcell{3}{7}{5.5}{7.5}{6}{1}{2}
\fxmsolid{3}{(7,5.5) -- (7.5,5.5) -- (7.5,6) -- (7,6) -- cycle}
\fxmcell{3}{7.5}{5.5}{8}{6}{1}{2}
\fxmsolid{3}{(7.5,5.5) -- (8,5.5) -- (8,6) -- (7.5,6) -- cycle}
\fxmopen{3}{(6.5,5) -- (7,5) -- (7,5.165022) -- (6.774767,5.5) -- (6.5,5.5) -- cycle}
\fxmopen{3}{(7,5) -- (7.22103,5) -- (7,5.238627) -- cycle}
\fxmdual{3}{6.75}{5}{7.25}{5.5}{1}
\fxmface{3}{2}{7.25}{5}{0.5}{l'}{0.47}
\fxmface{3}{2}{6.75}{5.5}{0.5}{k}{0.53}
\fxmface{3}{2}{7.25}{5.5}{0.5}{k'}{0.00}
\fxmface{3}{1}{7.5}{5.25}{0.5}{g'}{0.00}
\fxmface{3}{1}{6.5}{5.25}{0.5}{g}{1.00}
\fxmface{3}{1}{7}{5.25}{0.5}{f}{0.40}
\fxmface{3}{2}{6.75}{5}{0.5}{l}{1.00}
\fxmdualface{3}{6.75}{5}{6.75}{5.5}{\tfrac12(\aperture{f}\,f+\aperture{g}\,g)}{l}
\fxmdualface{3}{6.75}{5}{7.25}{5}{\tfrac12(\aperture{l}\,l+\aperture{l'}\,l')}{b}
\fxmdualface{3}{6.75}{5.5}{7.25}{5.5}{\tfrac12(\aperture{k}\,k+\aperture{k'}\,k')}{t}
\fxmdualface{3}{7.25}{5}{7.25}{5.5}{\tfrac12(\aperture{f}\,f+\aperture{g'}\,g')}{r}

      \node[tr centroid] at (5.5,5.5) {};
      \node[tr centroid] at (6.5,5.5) {};
    \end{tikzpicture}
    \caption{Uniform face.}
    \label{fig:momentum:dual:uniform}
  \end{subfigure}\\[6pt]
  \begin{subfigure}{\linewidth}
    \centering
    \begin{tikzpicture}[x=3.4cm,y=3.4cm]
      \begin{scope}
        \clip (6.38,4.86) rectangle (7.62,5.64);
        \fxmpasses{3}
      \end{scope}
      \fxmfacelabels{3}{0.015}{-0.02}{1}{-0.3}
      \fxmnone
      \def\fxmdualface#1#2#3#4#5#6#7{\ifnum#1=3
        \draw[trpurple,line width=1.1pt] (#2,#3) -- (#4,#5);\fi}%
      % Generated by docs/techref/figures/generate.jl; do not edit.
\fxmcell{1}{4}{4}{5}{5}{0}{0}
\fxmcell{1}{5}{4}{6}{5}{0}{0}
\fxmcell{1}{6}{4}{7}{5}{0}{0}
\fxmcell{1}{7}{4}{8}{5}{0}{0}
\fxmcell{1}{4}{5}{5}{6}{0}{0}
\fxmcell{1}{5}{5}{6}{6}{0}{0}
\fxmcell{1}{6}{5}{7}{6}{0}{0}
\fxmcell{1}{7}{5}{8}{6}{0}{0}
\fxmcell{1}{4}{6}{5}{7}{0}{0}
\fxmcell{1}{5}{6}{6}{7}{0}{0}
\fxmcell{1}{6}{6}{7}{7}{0}{0}
\fxmcell{1}{7}{6}{8}{7}{0}{0}
\fxmdual{1}{4.5}{4}{5.5}{5}{0}
\fxmdual{1}{4.5}{5}{5.5}{6}{0}
\fxmdual{1}{4.5}{6}{5.5}{7}{0}
\fxmdual{1}{5.5}{4}{6.5}{5}{0}
\fxmdual{1}{5.5}{5}{6.5}{6}{1}
\fxmdual{1}{5.5}{6}{6.5}{7}{0}
\fxmdual{1}{6.5}{4}{7.5}{5}{0}
\fxmdual{1}{6.5}{5}{7.5}{6}{0}
\fxmdual{1}{6.5}{6}{7.5}{7}{0}
\fxmface{1}{1}{6}{5.5}{1}{f}{1.00}
\fxmface{1}{2}{6.5}{6}{1}{k'}{1.00}
\fxmface{1}{2}{5.5}{5}{1}{l}{1.00}
\fxmface{1}{2}{6.5}{5}{1}{l'}{1.00}
\fxmface{1}{1}{7}{5.5}{1}{g'}{1.00}
\fxmface{1}{1}{5}{5.5}{1}{g}{1.00}
\fxmface{1}{2}{5.5}{6}{1}{k}{1.00}
\fxmdualface{1}{5.5}{5}{5.5}{6}{\tfrac{1}{2}\,f+\tfrac{1}{2}\,g}{l}
\fxmdualface{1}{6.5}{5}{6.5}{6}{\tfrac{1}{2}\,f+\tfrac{1}{2}\,g'}{r}
\fxmdualface{1}{5.5}{5}{6.5}{5}{\tfrac{1}{2}\,l+\tfrac{1}{2}\,l'}{b}
\fxmdualface{1}{5.5}{6}{6.5}{6}{\tfrac{1}{2}\,k+\tfrac{1}{2}\,k'}{t}
\fxmcell{2}{2}{3}{3}{4}{0}{0}
\fxmcell{2}{3}{3}{4}{4}{0}{0}
\fxmcell{2}{4}{3}{5}{4}{0}{0}
\fxmcell{2}{5}{3}{6}{4}{0}{0}
\fxmcell{2}{2}{4}{3}{5}{0}{0}
\fxmcell{2}{3}{4}{4}{5}{0}{0}
\fxmcell{2}{2}{5}{3}{6}{0}{0}
\fxmcell{2}{3}{5}{4}{6}{0}{0}
\fxmcell{2}{2}{6}{3}{7}{0}{0}
\fxmcell{2}{3}{6}{4}{7}{0}{0}
\fxmcell{2}{2}{7}{3}{8}{0}{0}
\fxmcell{2}{3}{7}{4}{8}{0}{0}
\fxmcell{2}{4}{4}{4.5}{4.5}{1}{0}
\fxmcell{2}{4.5}{4}{5}{4.5}{1}{0}
\fxmcell{2}{5}{4}{5.5}{4.5}{1}{0}
\fxmcell{2}{4}{4.5}{4.5}{5}{1}{0}
\fxmcell{2}{4.5}{4.5}{5}{5}{1}{0}
\fxmcell{2}{5}{4.5}{5.5}{5}{1}{0}
\fxmcell{2}{4}{5}{4.5}{5.5}{1}{0}
\fxmcell{2}{4.5}{5}{5}{5.5}{1}{0}
\fxmcell{2}{5}{5}{5.5}{5.5}{1}{0}
\fxmcell{2}{4}{5.5}{4.5}{6}{1}{0}
\fxmcell{2}{4.5}{5.5}{5}{6}{1}{0}
\fxmcell{2}{5}{5.5}{5.5}{6}{1}{0}
\fxmcell{2}{4}{6}{4.5}{6.5}{1}{0}
\fxmcell{2}{4.5}{6}{5}{6.5}{1}{0}
\fxmcell{2}{5}{6}{5.5}{6.5}{1}{0}
\fxmcell{2}{4}{6.5}{4.5}{7}{1}{0}
\fxmcell{2}{4.5}{6.5}{5}{7}{1}{0}
\fxmcell{2}{5}{6.5}{5.5}{7}{1}{0}
\fxmcell{2}{4}{7}{4.5}{7.5}{1}{0}
\fxmcell{2}{4.5}{7}{5}{7.5}{1}{0}
\fxmcell{2}{5}{7}{5.5}{7.5}{1}{0}
\fxmdual{2}{1.5}{3}{2.5}{4}{0}
\fxmdual{2}{1.5}{4}{2.5}{5}{0}
\fxmdual{2}{1.5}{5}{2.5}{6}{0}
\fxmdual{2}{1.5}{6}{2.5}{7}{0}
\fxmdual{2}{1.5}{7}{2.5}{8}{0}
\fxmdual{2}{2.5}{3}{3.5}{4}{0}
\fxmdual{2}{2.5}{4}{3.5}{5}{0}
\fxmdual{2}{2.5}{5}{3.5}{6}{0}
\fxmdual{2}{2.5}{6}{3.5}{7}{0}
\fxmdual{2}{2.5}{7}{3.5}{8}{0}
\fxmdual{2}{3.5}{3}{4.5}{4}{0}
\fxmdual{2}{3.5}{4}{4.25}{4.5}{0}
\fxmdual{2}{3.5}{4.5}{4.25}{5}{0}
\fxmdual{2}{3.5}{5}{4.25}{5.5}{1}
\fxmdual{2}{3.5}{5.5}{4.25}{6}{0}
\fxmdual{2}{3.5}{6}{4.25}{6.5}{0}
\fxmdual{2}{3.5}{6.5}{4.25}{7}{0}
\fxmdual{2}{3.5}{7}{4.25}{7.5}{0}
\fxmdual{2}{4.25}{4}{4.75}{4.5}{0}
\fxmdual{2}{4.25}{4.5}{4.75}{5}{0}
\fxmdual{2}{4.25}{5}{4.75}{5.5}{0}
\fxmdual{2}{4.25}{5.5}{4.75}{6}{0}
\fxmdual{2}{4.25}{6}{4.75}{6.5}{0}
\fxmdual{2}{4.25}{6.5}{4.75}{7}{0}
\fxmdual{2}{4.25}{7}{4.75}{7.5}{0}
\fxmdual{2}{4.5}{3}{5.5}{4}{0}
\fxmdual{2}{4.75}{4}{5.25}{4.5}{0}
\fxmdual{2}{4.75}{4.5}{5.25}{5}{0}
\fxmdual{2}{4.75}{5}{5.25}{5.5}{0}
\fxmdual{2}{4.75}{5.5}{5.25}{6}{0}
\fxmdual{2}{4.75}{6}{5.25}{6.5}{0}
\fxmdual{2}{4.75}{6.5}{5.25}{7}{0}
\fxmdual{2}{4.75}{7}{5.25}{7.5}{0}
\fxmface{2}{1}{4}{5.75}{0.5}{f_2}{1.00}
\fxmface{2}{1}{4}{5.25}{0.5}{f_1}{1.00}
\fxmface{2}{1}{4.5}{5.25}{0.5}{h_1}{1.00}
\fxmface{2}{2}{3.5}{6}{1}{\bar{k}^{+}}{1.00}
\fxmface{2}{2}{4.25}{5.5}{0.5}{k}{1.00}
\fxmface{2}{1}{3}{5.5}{1}{g}{1.00}
\fxmface{2}{2}{3.5}{5}{1}{\bar{k}^{-}}{1.00}
\fxmface{2}{2}{4.25}{5}{0.5}{k_1}{1.00}
\fxmlink{2}{4}{5.25}{3}{5.5}{1}{1}{\tfrac{1}{2}\,f_1+\tfrac{1}{4}\,g}{0}
\fxmlink{2}{4}{5.25}{4}{4.25}{2}{1}{\tfrac{1}{8}\,\bar{k}^{-}}{-1}
\fxmlink{2}{4}{5.25}{4}{4.75}{2}{1}{\tfrac{1}{8}\,\bar{k}^{-}+\tfrac{1}{2}\,k_1}{-1}
\fxmlink{2}{4}{5.25}{4}{5.75}{1}{1}{-\tfrac{1}{4}\,f_1+\tfrac{1}{4}\,f_2}{1}
\fxmlink{2}{4}{5.25}{4}{5.75}{2}{1}{\tfrac{1}{2}\,k}{-1}
\fxmlink{2}{4}{5.25}{4}{6.25}{2}{1}{\tfrac{1}{8}\,\bar{k}^{+}}{-1}
\fxmlink{2}{4}{5.25}{4}{6.75}{2}{1}{\tfrac{1}{8}\,\bar{k}^{+}}{-1}
\fxmlink{2}{4}{5.25}{4.5}{5.25}{1}{1}{\tfrac{1}{2}\,f_1+\tfrac{1}{2}\,h_1}{0}
\fxmhanging{2}{4}{4.5}
\fxmhanging{2}{4}{4.5}
\fxmhanging{2}{4}{5.5}
\fxmhanging{2}{4}{5.5}
\fxmhanging{2}{4}{6.5}
\fxmhanging{2}{4}{6.5}
\fxmcell{3}{6}{4.5}{6.5}{5}{1}{0}
\fxmcell{3}{6.5}{4.5}{7}{5}{1}{0}
\fxmcell{3}{7}{4.5}{7.5}{5}{1}{1}
\fxmsolid{3}{(7.5,4.894926) -- (7.5,5) -- (7.244029,5) -- cycle}
\fxmwall{3}{7.5}{4.894926}{7.244029}{5}
\fxmcell{3}{7.5}{4.5}{8}{5}{1}{1}
\fxmsolid{3}{(8,4.783878) -- (8,5) -- (7.5,5) -- (7.5,4.857327) -- cycle}
\fxmwall{3}{8}{4.783878}{7.5}{4.857327}
\fxmcell{3}{6}{5}{6.5}{5.5}{1}{0}
\fxmcell{3}{6.5}{5}{7}{5.5}{1}{1}
\fxmsolid{3}{(7,5.165022) -- (7,5.5) -- (6.774767,5.5) -- cycle}
\fxmwall{3}{7}{5.165022}{6.774767}{5.5}
\fxmcell{3}{7}{5}{7.5}{5.5}{1}{1}
\fxmsolid{3}{(7.22103,5) -- (7.5,5) -- (7.5,5.5) -- (7,5.5) -- (7,5.238627) -- cycle}
\fxmwall{3}{7.22103}{5}{7}{5.238627}
\fxmcell{3}{7.5}{5}{8}{5.5}{1}{2}
\fxmsolid{3}{(7.5,5) -- (8,5) -- (8,5.5) -- (7.5,5.5) -- cycle}
\fxmcell{3}{6}{5.5}{6.5}{6}{1}{0}
\fxmcell{3}{6.5}{5.5}{7}{6}{1}{1}
\fxmsolid{3}{(6.753218,5.5) -- (7,5.5) -- (7,6) -- (6.528991,6) -- cycle}
\fxmwall{3}{6.753218}{5.5}{6.528991}{6}
\fxmcell{3}{7}{5.5}{7.5}{6}{1}{2}
\fxmsolid{3}{(7,5.5) -- (7.5,5.5) -- (7.5,6) -- (7,6) -- cycle}
\fxmcell{3}{7.5}{5.5}{8}{6}{1}{2}
\fxmsolid{3}{(7.5,5.5) -- (8,5.5) -- (8,6) -- (7.5,6) -- cycle}
\fxmopen{3}{(6.5,5) -- (7,5) -- (7,5.165022) -- (6.774767,5.5) -- (6.5,5.5) -- cycle}
\fxmopen{3}{(7,5) -- (7.22103,5) -- (7,5.238627) -- cycle}
\fxmdual{3}{6.75}{5}{7.25}{5.5}{1}
\fxmface{3}{2}{7.25}{5}{0.5}{l'}{0.47}
\fxmface{3}{2}{6.75}{5.5}{0.5}{k}{0.53}
\fxmface{3}{2}{7.25}{5.5}{0.5}{k'}{0.00}
\fxmface{3}{1}{7.5}{5.25}{0.5}{g'}{0.00}
\fxmface{3}{1}{6.5}{5.25}{0.5}{g}{1.00}
\fxmface{3}{1}{7}{5.25}{0.5}{f}{0.40}
\fxmface{3}{2}{6.75}{5}{0.5}{l}{1.00}
\fxmdualface{3}{6.75}{5}{6.75}{5.5}{\tfrac12(\aperture{f}\,f+\aperture{g}\,g)}{l}
\fxmdualface{3}{6.75}{5}{7.25}{5}{\tfrac12(\aperture{l}\,l+\aperture{l'}\,l')}{b}
\fxmdualface{3}{6.75}{5.5}{7.25}{5.5}{\tfrac12(\aperture{k}\,k+\aperture{k'}\,k')}{t}
\fxmdualface{3}{7.25}{5}{7.25}{5.5}{\tfrac12(\aperture{f}\,f+\aperture{g'}\,g')}{r}

    \end{tikzpicture}
    \caption{Face between two cut cells.}
    \label{fig:momentum:dual:cut}
  \end{subfigure}
  \end{minipage}%
  \hfill
  \begin{subfigure}[b]{0.535\textwidth}
    \centering
    \begin{tikzpicture}[x=2.0cm,y=2.0cm]
      \begin{scope}
        \clip (2.6,4.1) rectangle (4.95,6.9);
        \fxmpasses{2}
      \end{scope}
      \fxmnone
      \def\fxmhanging#1#2#3{\ifnum#1=2 \node[tr vertex,minimum size=3.6pt] at (#2,#3) {};\fi}%
      % Generated by docs/techref/figures/generate.jl; do not edit.
\fxmcell{1}{4}{4}{5}{5}{0}{0}
\fxmcell{1}{5}{4}{6}{5}{0}{0}
\fxmcell{1}{6}{4}{7}{5}{0}{0}
\fxmcell{1}{7}{4}{8}{5}{0}{0}
\fxmcell{1}{4}{5}{5}{6}{0}{0}
\fxmcell{1}{5}{5}{6}{6}{0}{0}
\fxmcell{1}{6}{5}{7}{6}{0}{0}
\fxmcell{1}{7}{5}{8}{6}{0}{0}
\fxmcell{1}{4}{6}{5}{7}{0}{0}
\fxmcell{1}{5}{6}{6}{7}{0}{0}
\fxmcell{1}{6}{6}{7}{7}{0}{0}
\fxmcell{1}{7}{6}{8}{7}{0}{0}
\fxmdual{1}{4.5}{4}{5.5}{5}{0}
\fxmdual{1}{4.5}{5}{5.5}{6}{0}
\fxmdual{1}{4.5}{6}{5.5}{7}{0}
\fxmdual{1}{5.5}{4}{6.5}{5}{0}
\fxmdual{1}{5.5}{5}{6.5}{6}{1}
\fxmdual{1}{5.5}{6}{6.5}{7}{0}
\fxmdual{1}{6.5}{4}{7.5}{5}{0}
\fxmdual{1}{6.5}{5}{7.5}{6}{0}
\fxmdual{1}{6.5}{6}{7.5}{7}{0}
\fxmface{1}{1}{6}{5.5}{1}{f}{1.00}
\fxmface{1}{2}{6.5}{6}{1}{k'}{1.00}
\fxmface{1}{2}{5.5}{5}{1}{l}{1.00}
\fxmface{1}{2}{6.5}{5}{1}{l'}{1.00}
\fxmface{1}{1}{7}{5.5}{1}{g'}{1.00}
\fxmface{1}{1}{5}{5.5}{1}{g}{1.00}
\fxmface{1}{2}{5.5}{6}{1}{k}{1.00}
\fxmdualface{1}{5.5}{5}{5.5}{6}{\tfrac{1}{2}\,f+\tfrac{1}{2}\,g}{l}
\fxmdualface{1}{6.5}{5}{6.5}{6}{\tfrac{1}{2}\,f+\tfrac{1}{2}\,g'}{r}
\fxmdualface{1}{5.5}{5}{6.5}{5}{\tfrac{1}{2}\,l+\tfrac{1}{2}\,l'}{b}
\fxmdualface{1}{5.5}{6}{6.5}{6}{\tfrac{1}{2}\,k+\tfrac{1}{2}\,k'}{t}
\fxmcell{2}{2}{3}{3}{4}{0}{0}
\fxmcell{2}{3}{3}{4}{4}{0}{0}
\fxmcell{2}{4}{3}{5}{4}{0}{0}
\fxmcell{2}{5}{3}{6}{4}{0}{0}
\fxmcell{2}{2}{4}{3}{5}{0}{0}
\fxmcell{2}{3}{4}{4}{5}{0}{0}
\fxmcell{2}{2}{5}{3}{6}{0}{0}
\fxmcell{2}{3}{5}{4}{6}{0}{0}
\fxmcell{2}{2}{6}{3}{7}{0}{0}
\fxmcell{2}{3}{6}{4}{7}{0}{0}
\fxmcell{2}{2}{7}{3}{8}{0}{0}
\fxmcell{2}{3}{7}{4}{8}{0}{0}
\fxmcell{2}{4}{4}{4.5}{4.5}{1}{0}
\fxmcell{2}{4.5}{4}{5}{4.5}{1}{0}
\fxmcell{2}{5}{4}{5.5}{4.5}{1}{0}
\fxmcell{2}{4}{4.5}{4.5}{5}{1}{0}
\fxmcell{2}{4.5}{4.5}{5}{5}{1}{0}
\fxmcell{2}{5}{4.5}{5.5}{5}{1}{0}
\fxmcell{2}{4}{5}{4.5}{5.5}{1}{0}
\fxmcell{2}{4.5}{5}{5}{5.5}{1}{0}
\fxmcell{2}{5}{5}{5.5}{5.5}{1}{0}
\fxmcell{2}{4}{5.5}{4.5}{6}{1}{0}
\fxmcell{2}{4.5}{5.5}{5}{6}{1}{0}
\fxmcell{2}{5}{5.5}{5.5}{6}{1}{0}
\fxmcell{2}{4}{6}{4.5}{6.5}{1}{0}
\fxmcell{2}{4.5}{6}{5}{6.5}{1}{0}
\fxmcell{2}{5}{6}{5.5}{6.5}{1}{0}
\fxmcell{2}{4}{6.5}{4.5}{7}{1}{0}
\fxmcell{2}{4.5}{6.5}{5}{7}{1}{0}
\fxmcell{2}{5}{6.5}{5.5}{7}{1}{0}
\fxmcell{2}{4}{7}{4.5}{7.5}{1}{0}
\fxmcell{2}{4.5}{7}{5}{7.5}{1}{0}
\fxmcell{2}{5}{7}{5.5}{7.5}{1}{0}
\fxmdual{2}{1.5}{3}{2.5}{4}{0}
\fxmdual{2}{1.5}{4}{2.5}{5}{0}
\fxmdual{2}{1.5}{5}{2.5}{6}{0}
\fxmdual{2}{1.5}{6}{2.5}{7}{0}
\fxmdual{2}{1.5}{7}{2.5}{8}{0}
\fxmdual{2}{2.5}{3}{3.5}{4}{0}
\fxmdual{2}{2.5}{4}{3.5}{5}{0}
\fxmdual{2}{2.5}{5}{3.5}{6}{0}
\fxmdual{2}{2.5}{6}{3.5}{7}{0}
\fxmdual{2}{2.5}{7}{3.5}{8}{0}
\fxmdual{2}{3.5}{3}{4.5}{4}{0}
\fxmdual{2}{3.5}{4}{4.25}{4.5}{0}
\fxmdual{2}{3.5}{4.5}{4.25}{5}{0}
\fxmdual{2}{3.5}{5}{4.25}{5.5}{1}
\fxmdual{2}{3.5}{5.5}{4.25}{6}{0}
\fxmdual{2}{3.5}{6}{4.25}{6.5}{0}
\fxmdual{2}{3.5}{6.5}{4.25}{7}{0}
\fxmdual{2}{3.5}{7}{4.25}{7.5}{0}
\fxmdual{2}{4.25}{4}{4.75}{4.5}{0}
\fxmdual{2}{4.25}{4.5}{4.75}{5}{0}
\fxmdual{2}{4.25}{5}{4.75}{5.5}{0}
\fxmdual{2}{4.25}{5.5}{4.75}{6}{0}
\fxmdual{2}{4.25}{6}{4.75}{6.5}{0}
\fxmdual{2}{4.25}{6.5}{4.75}{7}{0}
\fxmdual{2}{4.25}{7}{4.75}{7.5}{0}
\fxmdual{2}{4.5}{3}{5.5}{4}{0}
\fxmdual{2}{4.75}{4}{5.25}{4.5}{0}
\fxmdual{2}{4.75}{4.5}{5.25}{5}{0}
\fxmdual{2}{4.75}{5}{5.25}{5.5}{0}
\fxmdual{2}{4.75}{5.5}{5.25}{6}{0}
\fxmdual{2}{4.75}{6}{5.25}{6.5}{0}
\fxmdual{2}{4.75}{6.5}{5.25}{7}{0}
\fxmdual{2}{4.75}{7}{5.25}{7.5}{0}
\fxmface{2}{1}{4}{5.75}{0.5}{f_2}{1.00}
\fxmface{2}{1}{4}{5.25}{0.5}{f_1}{1.00}
\fxmface{2}{1}{4.5}{5.25}{0.5}{h_1}{1.00}
\fxmface{2}{2}{3.5}{6}{1}{\bar{k}^{+}}{1.00}
\fxmface{2}{2}{4.25}{5.5}{0.5}{k}{1.00}
\fxmface{2}{1}{3}{5.5}{1}{g}{1.00}
\fxmface{2}{2}{3.5}{5}{1}{\bar{k}^{-}}{1.00}
\fxmface{2}{2}{4.25}{5}{0.5}{k_1}{1.00}
\fxmlink{2}{4}{5.25}{3}{5.5}{1}{1}{\tfrac{1}{2}\,f_1+\tfrac{1}{4}\,g}{0}
\fxmlink{2}{4}{5.25}{4}{4.25}{2}{1}{\tfrac{1}{8}\,\bar{k}^{-}}{-1}
\fxmlink{2}{4}{5.25}{4}{4.75}{2}{1}{\tfrac{1}{8}\,\bar{k}^{-}+\tfrac{1}{2}\,k_1}{-1}
\fxmlink{2}{4}{5.25}{4}{5.75}{1}{1}{-\tfrac{1}{4}\,f_1+\tfrac{1}{4}\,f_2}{1}
\fxmlink{2}{4}{5.25}{4}{5.75}{2}{1}{\tfrac{1}{2}\,k}{-1}
\fxmlink{2}{4}{5.25}{4}{6.25}{2}{1}{\tfrac{1}{8}\,\bar{k}^{+}}{-1}
\fxmlink{2}{4}{5.25}{4}{6.75}{2}{1}{\tfrac{1}{8}\,\bar{k}^{+}}{-1}
\fxmlink{2}{4}{5.25}{4.5}{5.25}{1}{1}{\tfrac{1}{2}\,f_1+\tfrac{1}{2}\,h_1}{0}
\fxmhanging{2}{4}{4.5}
\fxmhanging{2}{4}{4.5}
\fxmhanging{2}{4}{5.5}
\fxmhanging{2}{4}{5.5}
\fxmhanging{2}{4}{6.5}
\fxmhanging{2}{4}{6.5}
\fxmcell{3}{6}{4.5}{6.5}{5}{1}{0}
\fxmcell{3}{6.5}{4.5}{7}{5}{1}{0}
\fxmcell{3}{7}{4.5}{7.5}{5}{1}{1}
\fxmsolid{3}{(7.5,4.894926) -- (7.5,5) -- (7.244029,5) -- cycle}
\fxmwall{3}{7.5}{4.894926}{7.244029}{5}
\fxmcell{3}{7.5}{4.5}{8}{5}{1}{1}
\fxmsolid{3}{(8,4.783878) -- (8,5) -- (7.5,5) -- (7.5,4.857327) -- cycle}
\fxmwall{3}{8}{4.783878}{7.5}{4.857327}
\fxmcell{3}{6}{5}{6.5}{5.5}{1}{0}
\fxmcell{3}{6.5}{5}{7}{5.5}{1}{1}
\fxmsolid{3}{(7,5.165022) -- (7,5.5) -- (6.774767,5.5) -- cycle}
\fxmwall{3}{7}{5.165022}{6.774767}{5.5}
\fxmcell{3}{7}{5}{7.5}{5.5}{1}{1}
\fxmsolid{3}{(7.22103,5) -- (7.5,5) -- (7.5,5.5) -- (7,5.5) -- (7,5.238627) -- cycle}
\fxmwall{3}{7.22103}{5}{7}{5.238627}
\fxmcell{3}{7.5}{5}{8}{5.5}{1}{2}
\fxmsolid{3}{(7.5,5) -- (8,5) -- (8,5.5) -- (7.5,5.5) -- cycle}
\fxmcell{3}{6}{5.5}{6.5}{6}{1}{0}
\fxmcell{3}{6.5}{5.5}{7}{6}{1}{1}
\fxmsolid{3}{(6.753218,5.5) -- (7,5.5) -- (7,6) -- (6.528991,6) -- cycle}
\fxmwall{3}{6.753218}{5.5}{6.528991}{6}
\fxmcell{3}{7}{5.5}{7.5}{6}{1}{2}
\fxmsolid{3}{(7,5.5) -- (7.5,5.5) -- (7.5,6) -- (7,6) -- cycle}
\fxmcell{3}{7.5}{5.5}{8}{6}{1}{2}
\fxmsolid{3}{(7.5,5.5) -- (8,5.5) -- (8,6) -- (7.5,6) -- cycle}
\fxmopen{3}{(6.5,5) -- (7,5) -- (7,5.165022) -- (6.774767,5.5) -- (6.5,5.5) -- cycle}
\fxmopen{3}{(7,5) -- (7.22103,5) -- (7,5.238627) -- cycle}
\fxmdual{3}{6.75}{5}{7.25}{5.5}{1}
\fxmface{3}{2}{7.25}{5}{0.5}{l'}{0.47}
\fxmface{3}{2}{6.75}{5.5}{0.5}{k}{0.53}
\fxmface{3}{2}{7.25}{5.5}{0.5}{k'}{0.00}
\fxmface{3}{1}{7.5}{5.25}{0.5}{g'}{0.00}
\fxmface{3}{1}{6.5}{5.25}{0.5}{g}{1.00}
\fxmface{3}{1}{7}{5.25}{0.5}{f}{0.40}
\fxmface{3}{2}{6.75}{5}{0.5}{l}{1.00}
\fxmdualface{3}{6.75}{5}{6.75}{5.5}{\tfrac12(\aperture{f}\,f+\aperture{g}\,g)}{l}
\fxmdualface{3}{6.75}{5}{7.25}{5}{\tfrac12(\aperture{l}\,l+\aperture{l'}\,l')}{b}
\fxmdualface{3}{6.75}{5.5}{7.25}{5.5}{\tfrac12(\aperture{k}\,k+\aperture{k'}\,k')}{t}
\fxmdualface{3}{7.25}{5}{7.25}{5.5}{\tfrac12(\aperture{f}\,f+\aperture{g'}\,g')}{r}

      % f_1 is labelled right of the segment and f_2 above its middle, clear
      % of the dual faces that meet at the hanging vertex; the label of the
      % coarse face above moves left, clear of f_2.
      \def\fxmfone{f_1}\def\fxmftwo{f_2}\def\fxmkplus{\bar{k}^{+}}%
      \def\fxmlabeladjust#1{\def\fxmthis{#1}%
        \ifx\fxmthis\fxmfone\def\fxmdx{0.03}\fi
        \ifx\fxmthis\fxmftwo\def\fxmfrac{0.25}\fi
        \ifx\fxmthis\fxmkplus\def\fxmfrac{-0.15}\fi}%
      \fxmfacelabels{2}{-0.03}{-0.03}{0}{-0.25}
      \def\fxmlabeladjust#1{}%
      \fxmnone
      % Dual faces of f_1 as edges between the two unknowns they join
      % (dashed: non-regular), numbered; the weights are listed beside.
      \ifdefined\fxmlinkno\else\newcount\fxmlinkno\fi
      \global\fxmlinkno=0
      \tikzset{fxm bend 0/.style={}, fxm bend 1/.style={bend left=60},
        fxm bend -1/.style={bend right=40}, fxm bend 2/.style={bend right=40}}%
      \def\fxmlink#1#2#3#4#5#6#7#8#9{\ifnum#1=2
        \global\advance\fxmlinkno by 1
        \ifnum#7=1 \def\fxmstyle{tr link nr}\else\def\fxmstyle{tr link}\fi
        \draw[\fxmstyle] (#2,#3) to[fxm bend #9]
          node[midway,circle,fill=white,draw=trpurple,line width=0.4pt,solid,
            inner sep=0.4pt,minimum size=7pt,font=\tiny,text=trpurple!70!black] {\the\fxmlinkno} (#4,#5);
        \node[tr link label,anchor=west] at (5.0,{6.8-0.3*\the\fxmlinkno}) {\the\fxmlinkno:\ $#8$};\fi}%
      % Generated by docs/techref/figures/generate.jl; do not edit.
\fxmcell{1}{4}{4}{5}{5}{0}{0}
\fxmcell{1}{5}{4}{6}{5}{0}{0}
\fxmcell{1}{6}{4}{7}{5}{0}{0}
\fxmcell{1}{7}{4}{8}{5}{0}{0}
\fxmcell{1}{4}{5}{5}{6}{0}{0}
\fxmcell{1}{5}{5}{6}{6}{0}{0}
\fxmcell{1}{6}{5}{7}{6}{0}{0}
\fxmcell{1}{7}{5}{8}{6}{0}{0}
\fxmcell{1}{4}{6}{5}{7}{0}{0}
\fxmcell{1}{5}{6}{6}{7}{0}{0}
\fxmcell{1}{6}{6}{7}{7}{0}{0}
\fxmcell{1}{7}{6}{8}{7}{0}{0}
\fxmdual{1}{4.5}{4}{5.5}{5}{0}
\fxmdual{1}{4.5}{5}{5.5}{6}{0}
\fxmdual{1}{4.5}{6}{5.5}{7}{0}
\fxmdual{1}{5.5}{4}{6.5}{5}{0}
\fxmdual{1}{5.5}{5}{6.5}{6}{1}
\fxmdual{1}{5.5}{6}{6.5}{7}{0}
\fxmdual{1}{6.5}{4}{7.5}{5}{0}
\fxmdual{1}{6.5}{5}{7.5}{6}{0}
\fxmdual{1}{6.5}{6}{7.5}{7}{0}
\fxmface{1}{1}{6}{5.5}{1}{f}{1.00}
\fxmface{1}{2}{6.5}{6}{1}{k'}{1.00}
\fxmface{1}{2}{5.5}{5}{1}{l}{1.00}
\fxmface{1}{2}{6.5}{5}{1}{l'}{1.00}
\fxmface{1}{1}{7}{5.5}{1}{g'}{1.00}
\fxmface{1}{1}{5}{5.5}{1}{g}{1.00}
\fxmface{1}{2}{5.5}{6}{1}{k}{1.00}
\fxmdualface{1}{5.5}{5}{5.5}{6}{\tfrac{1}{2}\,f+\tfrac{1}{2}\,g}{l}
\fxmdualface{1}{6.5}{5}{6.5}{6}{\tfrac{1}{2}\,f+\tfrac{1}{2}\,g'}{r}
\fxmdualface{1}{5.5}{5}{6.5}{5}{\tfrac{1}{2}\,l+\tfrac{1}{2}\,l'}{b}
\fxmdualface{1}{5.5}{6}{6.5}{6}{\tfrac{1}{2}\,k+\tfrac{1}{2}\,k'}{t}
\fxmcell{2}{2}{3}{3}{4}{0}{0}
\fxmcell{2}{3}{3}{4}{4}{0}{0}
\fxmcell{2}{4}{3}{5}{4}{0}{0}
\fxmcell{2}{5}{3}{6}{4}{0}{0}
\fxmcell{2}{2}{4}{3}{5}{0}{0}
\fxmcell{2}{3}{4}{4}{5}{0}{0}
\fxmcell{2}{2}{5}{3}{6}{0}{0}
\fxmcell{2}{3}{5}{4}{6}{0}{0}
\fxmcell{2}{2}{6}{3}{7}{0}{0}
\fxmcell{2}{3}{6}{4}{7}{0}{0}
\fxmcell{2}{2}{7}{3}{8}{0}{0}
\fxmcell{2}{3}{7}{4}{8}{0}{0}
\fxmcell{2}{4}{4}{4.5}{4.5}{1}{0}
\fxmcell{2}{4.5}{4}{5}{4.5}{1}{0}
\fxmcell{2}{5}{4}{5.5}{4.5}{1}{0}
\fxmcell{2}{4}{4.5}{4.5}{5}{1}{0}
\fxmcell{2}{4.5}{4.5}{5}{5}{1}{0}
\fxmcell{2}{5}{4.5}{5.5}{5}{1}{0}
\fxmcell{2}{4}{5}{4.5}{5.5}{1}{0}
\fxmcell{2}{4.5}{5}{5}{5.5}{1}{0}
\fxmcell{2}{5}{5}{5.5}{5.5}{1}{0}
\fxmcell{2}{4}{5.5}{4.5}{6}{1}{0}
\fxmcell{2}{4.5}{5.5}{5}{6}{1}{0}
\fxmcell{2}{5}{5.5}{5.5}{6}{1}{0}
\fxmcell{2}{4}{6}{4.5}{6.5}{1}{0}
\fxmcell{2}{4.5}{6}{5}{6.5}{1}{0}
\fxmcell{2}{5}{6}{5.5}{6.5}{1}{0}
\fxmcell{2}{4}{6.5}{4.5}{7}{1}{0}
\fxmcell{2}{4.5}{6.5}{5}{7}{1}{0}
\fxmcell{2}{5}{6.5}{5.5}{7}{1}{0}
\fxmcell{2}{4}{7}{4.5}{7.5}{1}{0}
\fxmcell{2}{4.5}{7}{5}{7.5}{1}{0}
\fxmcell{2}{5}{7}{5.5}{7.5}{1}{0}
\fxmdual{2}{1.5}{3}{2.5}{4}{0}
\fxmdual{2}{1.5}{4}{2.5}{5}{0}
\fxmdual{2}{1.5}{5}{2.5}{6}{0}
\fxmdual{2}{1.5}{6}{2.5}{7}{0}
\fxmdual{2}{1.5}{7}{2.5}{8}{0}
\fxmdual{2}{2.5}{3}{3.5}{4}{0}
\fxmdual{2}{2.5}{4}{3.5}{5}{0}
\fxmdual{2}{2.5}{5}{3.5}{6}{0}
\fxmdual{2}{2.5}{6}{3.5}{7}{0}
\fxmdual{2}{2.5}{7}{3.5}{8}{0}
\fxmdual{2}{3.5}{3}{4.5}{4}{0}
\fxmdual{2}{3.5}{4}{4.25}{4.5}{0}
\fxmdual{2}{3.5}{4.5}{4.25}{5}{0}
\fxmdual{2}{3.5}{5}{4.25}{5.5}{1}
\fxmdual{2}{3.5}{5.5}{4.25}{6}{0}
\fxmdual{2}{3.5}{6}{4.25}{6.5}{0}
\fxmdual{2}{3.5}{6.5}{4.25}{7}{0}
\fxmdual{2}{3.5}{7}{4.25}{7.5}{0}
\fxmdual{2}{4.25}{4}{4.75}{4.5}{0}
\fxmdual{2}{4.25}{4.5}{4.75}{5}{0}
\fxmdual{2}{4.25}{5}{4.75}{5.5}{0}
\fxmdual{2}{4.25}{5.5}{4.75}{6}{0}
\fxmdual{2}{4.25}{6}{4.75}{6.5}{0}
\fxmdual{2}{4.25}{6.5}{4.75}{7}{0}
\fxmdual{2}{4.25}{7}{4.75}{7.5}{0}
\fxmdual{2}{4.5}{3}{5.5}{4}{0}
\fxmdual{2}{4.75}{4}{5.25}{4.5}{0}
\fxmdual{2}{4.75}{4.5}{5.25}{5}{0}
\fxmdual{2}{4.75}{5}{5.25}{5.5}{0}
\fxmdual{2}{4.75}{5.5}{5.25}{6}{0}
\fxmdual{2}{4.75}{6}{5.25}{6.5}{0}
\fxmdual{2}{4.75}{6.5}{5.25}{7}{0}
\fxmdual{2}{4.75}{7}{5.25}{7.5}{0}
\fxmface{2}{1}{4}{5.75}{0.5}{f_2}{1.00}
\fxmface{2}{1}{4}{5.25}{0.5}{f_1}{1.00}
\fxmface{2}{1}{4.5}{5.25}{0.5}{h_1}{1.00}
\fxmface{2}{2}{3.5}{6}{1}{\bar{k}^{+}}{1.00}
\fxmface{2}{2}{4.25}{5.5}{0.5}{k}{1.00}
\fxmface{2}{1}{3}{5.5}{1}{g}{1.00}
\fxmface{2}{2}{3.5}{5}{1}{\bar{k}^{-}}{1.00}
\fxmface{2}{2}{4.25}{5}{0.5}{k_1}{1.00}
\fxmlink{2}{4}{5.25}{3}{5.5}{1}{1}{\tfrac{1}{2}\,f_1+\tfrac{1}{4}\,g}{0}
\fxmlink{2}{4}{5.25}{4}{4.25}{2}{1}{\tfrac{1}{8}\,\bar{k}^{-}}{-1}
\fxmlink{2}{4}{5.25}{4}{4.75}{2}{1}{\tfrac{1}{8}\,\bar{k}^{-}+\tfrac{1}{2}\,k_1}{-1}
\fxmlink{2}{4}{5.25}{4}{5.75}{1}{1}{-\tfrac{1}{4}\,f_1+\tfrac{1}{4}\,f_2}{1}
\fxmlink{2}{4}{5.25}{4}{5.75}{2}{1}{\tfrac{1}{2}\,k}{-1}
\fxmlink{2}{4}{5.25}{4}{6.25}{2}{1}{\tfrac{1}{8}\,\bar{k}^{+}}{-1}
\fxmlink{2}{4}{5.25}{4}{6.75}{2}{1}{\tfrac{1}{8}\,\bar{k}^{+}}{-1}
\fxmlink{2}{4}{5.25}{4.5}{5.25}{1}{1}{\tfrac{1}{2}\,f_1+\tfrac{1}{2}\,h_1}{0}
\fxmhanging{2}{4}{4.5}
\fxmhanging{2}{4}{4.5}
\fxmhanging{2}{4}{5.5}
\fxmhanging{2}{4}{5.5}
\fxmhanging{2}{4}{6.5}
\fxmhanging{2}{4}{6.5}
\fxmcell{3}{6}{4.5}{6.5}{5}{1}{0}
\fxmcell{3}{6.5}{4.5}{7}{5}{1}{0}
\fxmcell{3}{7}{4.5}{7.5}{5}{1}{1}
\fxmsolid{3}{(7.5,4.894926) -- (7.5,5) -- (7.244029,5) -- cycle}
\fxmwall{3}{7.5}{4.894926}{7.244029}{5}
\fxmcell{3}{7.5}{4.5}{8}{5}{1}{1}
\fxmsolid{3}{(8,4.783878) -- (8,5) -- (7.5,5) -- (7.5,4.857327) -- cycle}
\fxmwall{3}{8}{4.783878}{7.5}{4.857327}
\fxmcell{3}{6}{5}{6.5}{5.5}{1}{0}
\fxmcell{3}{6.5}{5}{7}{5.5}{1}{1}
\fxmsolid{3}{(7,5.165022) -- (7,5.5) -- (6.774767,5.5) -- cycle}
\fxmwall{3}{7}{5.165022}{6.774767}{5.5}
\fxmcell{3}{7}{5}{7.5}{5.5}{1}{1}
\fxmsolid{3}{(7.22103,5) -- (7.5,5) -- (7.5,5.5) -- (7,5.5) -- (7,5.238627) -- cycle}
\fxmwall{3}{7.22103}{5}{7}{5.238627}
\fxmcell{3}{7.5}{5}{8}{5.5}{1}{2}
\fxmsolid{3}{(7.5,5) -- (8,5) -- (8,5.5) -- (7.5,5.5) -- cycle}
\fxmcell{3}{6}{5.5}{6.5}{6}{1}{0}
\fxmcell{3}{6.5}{5.5}{7}{6}{1}{1}
\fxmsolid{3}{(6.753218,5.5) -- (7,5.5) -- (7,6) -- (6.528991,6) -- cycle}
\fxmwall{3}{6.753218}{5.5}{6.528991}{6}
\fxmcell{3}{7}{5.5}{7.5}{6}{1}{2}
\fxmsolid{3}{(7,5.5) -- (7.5,5.5) -- (7.5,6) -- (7,6) -- cycle}
\fxmcell{3}{7.5}{5.5}{8}{6}{1}{2}
\fxmsolid{3}{(7.5,5.5) -- (8,5.5) -- (8,6) -- (7.5,6) -- cycle}
\fxmopen{3}{(6.5,5) -- (7,5) -- (7,5.165022) -- (6.774767,5.5) -- (6.5,5.5) -- cycle}
\fxmopen{3}{(7,5) -- (7.22103,5) -- (7,5.238627) -- cycle}
\fxmdual{3}{6.75}{5}{7.25}{5.5}{1}
\fxmface{3}{2}{7.25}{5}{0.5}{l'}{0.47}
\fxmface{3}{2}{6.75}{5.5}{0.5}{k}{0.53}
\fxmface{3}{2}{7.25}{5.5}{0.5}{k'}{0.00}
\fxmface{3}{1}{7.5}{5.25}{0.5}{g'}{0.00}
\fxmface{3}{1}{6.5}{5.25}{0.5}{g}{1.00}
\fxmface{3}{1}{7}{5.25}{0.5}{f}{0.40}
\fxmface{3}{2}{6.75}{5}{0.5}{l}{1.00}
\fxmdualface{3}{6.75}{5}{6.75}{5.5}{\tfrac12(\aperture{f}\,f+\aperture{g}\,g)}{l}
\fxmdualface{3}{6.75}{5}{7.25}{5}{\tfrac12(\aperture{l}\,l+\aperture{l'}\,l')}{b}
\fxmdualface{3}{6.75}{5.5}{7.25}{5.5}{\tfrac12(\aperture{k}\,k+\aperture{k'}\,k')}{t}
\fxmdualface{3}{7.25}{5}{7.25}{5.5}{\tfrac12(\aperture{f}\,f+\aperture{g'}\,g')}{r}

      \node[tr small,fill=white,fill opacity=0.8,text opacity=1,inner sep=1pt] at (2.85,5.85) {$c$};
      \node[tr small,fill=white,fill opacity=0.8,text opacity=1,inner sep=1pt] at (4.8,5.4) {$c_1$};
    \end{tikzpicture}
    \caption{Fine segment $f_1$ of a hanging side.}
    \label{fig:momentum:dual:coarse-fine}
  \end{subfigure}
  \caption{Staggered control volumes of the component $\velo$ on faces
    $f\in\facesaxis{1}$ (dashed; green: the one of the face considered) and
    their dual faces (purple) with the weights $\dweight{g}{f'}$ of the
    primal fluxes $\meas{f'}\volumeflux_{f'}$, read from the code's schedules
    (labels: weight times face name; red: faces of $\facesaxis{1}$, blue:
    $\facesaxis{2}$). (a)~Uniform periodic mesh: each dual face averages two
    primal fluxes~\cref{eq:momentum:uniform-weights}. (b)~Disc fixture (cells
    of width $\half$, solid hatched): the weights are $\half\aperture{f'}$,
    labels give $\aperture{f'}$; green: the open parts of the two strips, half
    of whose area is $\opendualvol{f}=\chkMomCutOpenDual$ against
    $\dualvol{f}=\chkMomCutUncutDual$. (c)~Composite mesh of
    \cref{fig:vof:donating:coarse-fine}: all \chkMomCoarseLinks\ dual faces of the
    fine segment $f_1$, drawn as edges between the unknowns they join
    (dashed: non-regular; arcs join unknowns on the same line, bent right for
    $\facesaxis{1}$-fluxes and left for $\facesaxis{2}$-fluxes; open circles:
    hanging vertices); $\stagvol{f_1}=\chkMomCoarseVolume$.}
  \label{fig:momentum:dual}
\end{figure}
